\documentclass[12pt,openany]{book}
\usepackage{amsmath}
\usepackage{amssymb}
\usepackage{amsthm}
\usepackage{mathtools}
\usepackage[hidelinks]{hyperref}
\usepackage{booktabs}
\usepackage{geometry}
\usepackage{fancyhdr}
\usepackage{setspace}
\usepackage{microtype}
\usepackage{xcolor}
\usepackage{enumitem}
\usepackage{titlesec}
\usepackage{tocloft}
\usepackage{array}
\usepackage{tabularx}
\usepackage{longtable}
\usepackage{float}
\usepackage{caption}
\usepackage{multirow}
\usepackage{makecell}
\usepackage[numbers,sort&compress]{natbib} % used for bibliography management

% ============================================================
% GEOMETRY AND LAYOUT SETTINGS
% ============================================================
\geometry{margin=1in, includefoot}
\microtypesetup{expansion=true, final}
\singlespacing
\setlength{\parindent}{0.5in}
\setlength{\parskip}{0.5ex plus 0.2ex minus 0.1ex}
\setlength{\emergencystretch}{1.5em}
\tolerance=9999
\hbadness=10000
\vbadness=10000
\hyphenpenalty=10000
\exhyphenpenalty=10000
\raggedbottom
\sloppy

% ============================================================
% SECTION FORMATTING
% ============================================================
\titleformat{\chapter}[display]
    {\normalfont\Huge\bfseries}
    {\chaptertitlename\ \thechapter}{20pt}{\Huge}
\titlespacing*{\chapter}{0pt}{50pt}{40pt}
\titleformat{\section}{\normalfont\Large\bfseries}{\thesection}{1em}{}
\titlespacing*{\section}{0pt}{3.5ex plus 1ex minus 0.2ex}{2.3ex plus 0.2ex}
\titleformat{\subsection}{\normalfont\large\bfseries}{\thesubsection}{1em}{}
\titlespacing*{\subsection}{0pt}{3.25ex plus 1ex minus 0.2ex}{1.5ex plus 0.2ex}
\titleformat{\subsubsection}{\normalfont\normalsize\bfseries}{\thesubsubsection}{1em}{}
\titlespacing*{\subsubsection}{0pt}{3ex plus 0.5ex minus 0.2ex}{1.5ex plus 0.2ex}

% ============================================================
% TABLE OF CONTENTS
% ============================================================
\renewcommand{\cftchapleader}{\cftdotfill{\cftdotsep}}
\setlength{\cftchapindent}{0em}
\setlength{\cftsecindent}{1.5em}

% ============================================================
% HEADER AND FOOTER
% ============================================================
\pagestyle{fancy}
\fancyhf{}
\fancyhead[LE]{\thepage}
\fancyhead[RO]{\thepage}
\fancyhead[RE]{\nouppercase{\leftmark}}
\fancyhead[LO]{\nouppercase{\rightmark}}
\renewcommand{\headrulewidth}{0.4pt}
\renewcommand{\footrulewidth}{0pt}
\setlength{\headheight}{14.5pt}

% ============================================================
% EQUATION SETTINGS
% ============================================================
\allowdisplaybreaks[3]
\setlength{\abovedisplayskip}{6pt plus 2pt minus 3pt}
\setlength{\belowdisplayskip}{6pt plus 2pt minus 3pt}
\setlength{\abovedisplayshortskip}{3pt plus 1pt minus 1pt}
\setlength{\belowdisplayshortskip}{3pt plus 1pt minus 1pt}
\numberwithin{equation}{chapter}

% ============================================================
% THEOREM ENVIRONMENTS
% ============================================================
\theoremstyle{plain}
\newtheorem{theorem}{Theorem}[chapter]
\newtheorem{lemma}[theorem]{Lemma}
\newtheorem{proposition}[theorem]{Proposition}
\theoremstyle{definition}
\newtheorem{definition}[theorem]{Definition}
\newtheorem{principle}[theorem]{Principle}
\theoremstyle{remark}
\newtheorem{remark}[theorem]{Remark}

% ============================================================
% MISC
% ============================================================
\clubpenalty=10000
\widowpenalty=10000
\displaywidowpenalty=10000
\renewcommand{\arraystretch}{1.3}
\setlength{\tabcolsep}{5pt}
\setlist[enumerate]{leftmargin=*, partopsep=0.1ex}

% ============================================================
% CUSTOM MACROS
% ============================================================
\newcommand{\BigO}[1]{O\!\left(#1\right)}
\DeclareMathOperator{\CVaR}{CVaR}
\DeclareMathOperator{\VaR}{VaR}
\DeclareMathOperator{\Sharpe}{Sharpe}
\DeclareMathOperator{\tr}{tr}

% ============================================================
% TITLE
% ============================================================
\title{%
    \textbf{High-Level Architecture Whitepaper}\\[0.5em]
    \large A Multi-Agent, Cross-Market Quantitative Intelligence Platform\\
    for Fintech, Portfolio Risk, and Algorithmic Trading
}
\author{Hadrian Hu}
\date{\today}

\begin{document}

\maketitle
\tableofcontents

% ============================================================
% ABSTRACT
% ============================================================
\chapter*{Abstract}
\addcontentsline{toc}{chapter}{Abstract}

This whitepaper presents the complete high-level architecture of a
multi-agent, cross-market quantitative intelligence platform designed
for fintech research, portfolio risk management, market-regime classification,
hedging optimisation, and algorithmic trading. The platform is composed
of typed, composable software \emph{Operators} and \emph{Agents} implemented
across Python, C++, Rust, and Go, forming a coherent pipeline that transforms
raw market data into fully provenance-tracked, risk-gated trading decisions.

The architecture is structured around four financial domains: Portfolio and
Income (Earth), Hedging and Risk Control (Air), Volatility and Options Alpha
(Fire), and Macro and Cross-Market Dynamics (Water). These domains interact
through a shared operator graph that enforces temporal correctness, uncertainty
quantification, and risk-before-execution discipline at every stage.

The whitepaper establishes the foundational mathematical framework governing
each major subsystem, provides formal definitions of the operator composition
algebra, documents the data pipeline from raw ingestion through feature
extraction and regime inference, and specifies the testing taxonomy and
deployment constraints that distinguish this platform from naive backtesting
environments. The platform is explicitly designed for paper trading and
historical simulation as its initial mode of operation; live execution is
architecturally separated and requires deliberate, explicit enablement.

\paragraph*{Keywords:} quantitative finance, multi-agent systems, operator
composition, market regime, Bayesian inference, portfolio optimisation,
risk management, hedging, Kelly criterion, CVaR, temporal leakage, backtesting,
walk-forward validation, polyglot architecture, Python, C++, Rust, Go.

% ============================================================
% EXECUTIVE SUMMARY
% ============================================================
\chapter*{Executive Summary}
\addcontentsline{toc}{chapter}{Executive Summary}

\subparagraph{Platform Vision.}
The platform described in this whitepaper is a cross-market adaptive financial
decision system designed to operate coherently across equities, futures,
foreign exchange, and cryptocurrency venues without requiring fundamental
re-engineering at the strategy layer. Rather than deploying a single
monolithic predictive model that attempts to encode all market knowledge into
one undifferentiated parameter space, the platform composes small,
independently verifiable functional units called \emph{Operators} into larger
deliberative entities called \emph{Agents}, which are in turn orchestrated by
a top-level \emph{System} layer that governs capital allocation, risk
aggregation, and execution sequencing. Every decision the platform produces is
traceable to a deterministic chain of typed inputs, versioned model artifacts,
and immutable data snapshots, making every output auditable and fully
reproducible from first principles without reliance on global mutable state or
external side-effects. This compositional, provenance-first design is not an
engineering preference; it is a direct response to the observation that opaque,
monolithic financial models fail in ways that are difficult to diagnose,
impossible to certify, and catastrophic when deployed at scale.

\subparagraph{Core Problem Statement.}
Financial markets present a fundamentally ill-posed prediction problem that
cannot be resolved by simply increasing model capacity or training data volume.
Prices are driven simultaneously by the interaction of rational utility-
maximising agents, systematic behavioural biases such as momentum-chasing and
loss-aversion anchoring, low-frequency macroeconomic shocks originating in
central bank policy and fiscal intervention, medium-frequency regulatory and
structural changes including circuit-breaker rules and tick-size regimes, and
high-frequency emergent systemic dynamics arising from the co-location of
algorithmic market-makers operating on microsecond latencies. Each of these
forces operates on a different temporal scale, responds to different
information sets, and interacts non-linearly with the others, producing a joint
data-generating process whose stationarity properties are time-varying by
construction. No single model class---whether linear factor models, gradient-
boosted trees, recurrent neural networks, or transformer-based sequence
models---can capture all of these forces simultaneously with consistent,
out-of-sample accuracy across all market regimes and instrument classes.
The architecture therefore does not attempt to maximise a single predicted
profit metric, nor does it seek a unified ``alpha model'' that encodes
universal market knowledge. Instead, it poses the decision problem as a
structured sequence of conditional queries: \emph{Given all information
legitimately available at time~$t$, what market regime is most probable and
with what posterior uncertainty; what risk factors currently dominate the
cross-sectional covariance matrix; how confident is the ensemble of models in
its current assessment; and what portfolio allocation, directional hedge, or
discrete trading action provides the best risk-adjusted response given the
full distribution of plausible outcomes rather than merely the modal
forecast?} This reframing converts an intractable prediction problem into a
tractable decision problem under uncertainty, which admits principled
engineering solutions.

\subparagraph{Architectural Philosophy.}
The platform rests on six non-negotiable architectural principles that were
derived not from theoretical preference but from the systematic failure modes
observed in production financial systems and academic backtesting literature.
Each principle carries an associated enforcement mechanism that is embedded in
the type system, the data pipeline, or the deployment contract, rather than
left to developer discipline alone.
\begin{enumerate}
    \item \textbf{Operator composition over monolithic models.} Complex
          decision-making behaviour emerges from the composition of small,
          typed, independently testable, and formally specified functional
          units called Operators. Each Operator declares its input schema,
          output schema, temporal alignment requirements, and computational
          complexity class. Operators may be stateless transformations, such
          as a volatility normalisation step, or stateful accumulators, such
          as a hidden Markov model over rolling regime windows, but in all
          cases their dependencies are explicit and their outputs are
          inspectable. The composition graph is a directed acyclic graph
          whose topological sort determines evaluation order, enabling
          parallel evaluation of independent branches and deterministic
          replay of any historical computation path.

    \item \textbf{Risk gates every signal.} A signal that appears profitable
          under historical simulation never automatically produces a live or
          paper order. Every candidate action---whether a full position entry,
          a partial size adjustment, or a hedge overlay---must pass through a
          dedicated \texttt{RiskOperator} stage that evaluates the proposed
          action against a configurable policy set encompassing position
          concentration limits, drawdown thresholds, volatility-adjusted
          Kelly fraction bounds, correlation-based portfolio heat metrics, and
          regime-conditional exposure rules. The \texttt{RiskOperator} may
          reject the action outright, reduce its notional size, substitute an
          alternative instrument with lower correlation to existing holdings,
          delay execution pending a volatility cooldown window, or attach a
          mandatory hedge leg. Only actions that satisfy all active risk
          policies are forwarded to the execution layer; the decision and its
          policy evaluation trace are recorded unconditionally.

    \item \textbf{Simulation before live trading.} The platform is designed
          as a research and paper-trading environment first, and a live
          execution system second. Historical replay, walk-forward
          backtesting, Monte Carlo scenario generation, and stress testing
          against synthetic crisis regimes are first-class capabilities with
          dedicated pipeline stages and output schemas. The simulation
          environment enforces temporal correctness at the API boundary,
          making look-ahead access to future data a type error rather than a
          runtime failure. Paper execution is implemented as a high-fidelity
          market microstructure simulation that models bid-ask spread,
          partial fill probability, market impact as a function of order
          size relative to average daily volume, and queue position dynamics
          for limit orders, providing a realistic bridge between backtest
          performance and live execution expectation.

    \item \textbf{Provenance is first-class.} Every decision the platform
          produces is reconstructable from its constituent data sources,
          operator version identifiers, model checkpoint hashes, and
          configuration snapshots. The provenance record is not an optional
          audit log appended after the fact; it is a typed output of the
          execution pipeline that is produced unconditionally and stored
          atomically with the decision itself. This enables not only
          post-hoc audit and regulatory review but also automated regression
          testing: any change to an Operator, model, or schema version that
          produces a different decision for an identical historical input is
          flagged as a provenance divergence requiring explicit review before
          deployment.

    \item \textbf{Temporal correctness is a hard requirement.} A strategy
          that accesses data from time $t' > t$ during a decision that is
          nominally made at time $t$ is invalid regardless of its reported
          backtest profitability, Sharpe ratio, or maximum drawdown
          statistics. Temporal leakage---the inadvertent use of future
          information through lookahead bias, survivorship bias, or
          point-in-time data misalignment---is the single most common source
          of inflated backtest performance in the academic and practitioner
          literature. The platform enforces temporal correctness through a
          time-indexed data access interface that makes out-of-window queries
          a checked runtime error in simulation mode and a hard type error in
          the static analysis pipeline, ensuring that no operator or model
          can silently consume future data.

    \item \textbf{Uncertainty is information.} Disagreement between models,
          widening confidence intervals, and elevated entropy in regime
          probability distributions are not failure states to be suppressed
          or averaged away; they are first-class signals that carry decision-
          relevant information about the reliability of the current forecast.
          When the ensemble of regime classifiers disagrees significantly,
          the appropriate response is not to take the majority vote and
          proceed; it is to widen risk limits, reduce position sizes, increase
          hedge ratios, and surface the disagreement explicitly to any
          downstream consumer. Uncertainty quantification is therefore a
          mandatory output of every predictive Operator, expressed as a
          calibrated probability distribution or, where distributional
          assumptions are untenable, as a set-valued or interval-valued
          estimate that is propagated faithfully through the composition graph.
\end{enumerate}

\subparagraph{Scope.}
The initial platform implementation targets five primary operational modes:
historical replay for deterministic reconstruction of past market conditions;
walk-forward backtesting with configurable retraining and re-optimisation
windows to measure strategy robustness across distributional shifts; Monte
Carlo simulation for generating synthetic scenario ensembles that stress-test
strategy behaviour under tail-risk regimes not present in the historical
record; stress testing against hand-constructed crisis scenarios calibrated to
historical episodes such as the 2008 credit crisis, the 2020 pandemic shock,
and the 2022 rate-normalisation regime; and paper execution through a
microstructure-realistic simulation layer that records order lifecycle events
against a live or replayed market feed without committing capital. Live-trading
adapters that connect the execution layer to exchange APIs, prime broker
order-management systems, or direct-market-access gateways are architecturally
separated into a dedicated integration boundary with a well-specified contract,
and are deliberately absent from the initial implementation scope in order to
ensure that the research, simulation, and risk infrastructure reaches full
maturity and passes independent correctness review before any capital is placed
at risk.

% ============================================================
% KEY FINDINGS
% ============================================================
\chapter*{Key Findings}
\addcontentsline{toc}{chapter}{Key Findings}

\begin{table}[H]
\centering
\caption{Key architectural findings and their implications.}
\label{tab:key_findings}
\small
\begin{tabularx}{\textwidth}{p{0.5cm} p{5.5cm} p{8.5cm}}
\toprule
\textbf{\#} & \textbf{Finding} & \textbf{Implication} \\
\midrule
1 & Agents defined as Markdown files cannot be executable. &
    All agent capabilities must be typed source-code objects. \\
\midrule
2 & Candlestick patterns are features, not oracles. &
    Feature operators compute pattern presence; predictive value
    is determined empirically. \\
\midrule
3 & Market regime is a probability distribution, not a label. &
    \texttt{MarketRegimeState} carries a probability vector,
    confidence score, and uncertainty estimate. \\
\midrule
4 & Model disagreement is itself information. &
    Disagreement metrics must be computed and surfaced, not silently
    collapsed into a single decision. \\
\midrule
5 & Temporal leakage is a correctness failure. &
    A backtest that sees future data is invalid regardless of
    reported profitability. \\
\midrule
6 & Risk approval is a separate pipeline stage. &
    \texttt{RiskOperator} may reject, modify, size, delay, or hedge
    any proposed action. \\
\midrule
7 & Four languages require four defensible roles. &
    Python = research; C++ = numerics; Rust = integrity; Go = services. \\
\midrule
8 & Shared schemas prevent polyglot divergence. &
    A single language-neutral schema layer governs all shared types. \\
\midrule
9 & Configuration must not become a programming language. &
    Executable logic belongs in source code; configuration is
    declarative only. \\
\midrule
10 & Success is not maximum profit. &
    Success is defensible decisions under model disagreement, regime
    change, forecast failure, and elevated uncertainty. \\
\bottomrule
\end{tabularx}
\end{table}

% ============================================================
% CHAPTER 1: INTRODUCTION
% ============================================================
\chapter{Introduction}

\section{Motivation}

Modern financial markets generate terabytes of structured and unstructured
data every trading day: price ticks, order-book snapshots, options chains,
macroeconomic releases, central bank communications, earnings announcements,
news headlines, and satellite imagery of industrial facilities. Extracting
actionable, risk-adjusted decisions from this information requires a
principled engineering architecture, not merely a collection of prediction
scripts.

The dominant failure mode in quantitative finance is not model inaccuracy
per se. Backtesting frameworks routinely produce strategies with Sharpe
ratios above three that collapse to near zero in live trading. The root
causes are well documented: temporal data leakage, survivorship bias,
overfitting to in-sample regime characteristics, failure to account for
transaction costs and market impact, and the absence of robust risk
controls that would have prevented catastrophic drawdowns. This platform
is designed from first principles to address each of these failure modes
systematically.

\section{Scope and Objectives}

The platform described in this whitepaper has the following primary objectives:

\begin{enumerate}
    \item Provide a modular, composable operator framework that enables rapid
          development and independent testing of quantitative strategies.
    \item Enforce temporal correctness at the architectural level, making
          data leakage structurally impossible rather than merely discouraged
          by convention.
    \item Implement a multi-model market regime engine that produces calibrated
          probability distributions over regime states, with explicit uncertainty
          quantification.
    \item Gate every trading decision through a typed risk evaluation pipeline
          before it reaches an execution layer.
    \item Produce fully provenance-tracked outputs that allow any decision to
          be reconstructed and audited from its raw inputs.
    \item Support historical simulation, walk-forward backtesting, Monte Carlo
          analysis, and paper trading as first-class operational modes.
\end{enumerate}

The first objective—modularity and composability—addresses a pervasive structural deficiency in proprietary trading systems: the tendency for research codebases to accumulate tightly coupled, monolithic scripts in which data ingestion, feature engineering, signal generation, position sizing, and execution are interleaved without clear boundaries. When every component depends implicitly on the global state of every other component, the cost of replacing or improving any single element grows super-linearly with system complexity. The operator framework described in Chapter~3 resolves this by defining a strict functional interface for every computational unit: an operator consumes a temporally indexed, schema-validated input bundle and emits an output bundle of declared type, with no permitted side-effects on shared mutable state. Composition is achieved through a directed acyclic graph (DAG) of operator nodes, where edges carry typed data streams. This design enables any sub-graph to be extracted, unit-tested in isolation, and replaced without disturbing the remainder of the pipeline. The composability property further allows complex multi-asset, multi-horizon strategies to be expressed as hierarchical operator trees, with higher-level \emph{agent} nodes delegating sub-computations to specialised child operators while aggregating their outputs into coherent portfolio-level decisions.

The second objective—temporal correctness—is elevated from a software quality concern to an architectural guarantee. In a naive implementation, the burden of avoiding look-ahead bias falls entirely on the discipline of individual researchers, a burden that historical experience demonstrates is routinely violated under deadline pressure or through subtle misunderstanding of dataset construction. This platform instead encodes the information horizon of every data element as a first-class type parameter: every tensor, frame, or signal carries a \texttt{ValidFrom} timestamp indicating the earliest point in time at which that datum was legitimately observable by a market participant. The operator execution engine enforces, as a hard invariant, that no operator node may consume any input whose \texttt{ValidFrom} exceeds the operator's declared evaluation timestamp. Violations raise typed exceptions that propagate to the DAG scheduler and abort execution rather than silently contaminating downstream computations. This invariant is preserved across all operational modes—backtesting, walk-forward analysis, paper trading, and live execution—through a unified time-travel abstraction layer detailed in Chapter~4.

The third objective—calibrated probabilistic regime classification—reflects the empirical observation that a single statistical model of market dynamics is inadequate across the full range of macroeconomic and microstructure regimes that characterise modern equity, fixed-income, and derivatives markets. High-volatility stress regimes, low-volatility trending regimes, mean-reverting range-bound regimes, and crisis liquidity regimes each exhibit qualitatively different return distributions, autocorrelation structures, and cross-asset correlation patterns. The regime engine described in Chapter~5 maintains an ensemble of $K$ heterogeneous classifiers—including hidden Markov models, threshold autoregressive models, and gradient-boosted discriminative classifiers—each of which produces a posterior distribution $p_k(\mathcal{R}_t \mid \mathcal{F}_t)$ over a shared discrete regime vocabulary $\mathcal{R} = \{r_1, \ldots, r_M\}$ conditioned on the information set $\mathcal{F}_t$ available at time $t$. These posteriors are aggregated into a consensus distribution through a Bayesian model averaging scheme with dynamically updated model weights, yielding a final regime distribution $\bar{p}(\mathcal{R}_t \mid \mathcal{F}_t)$ that explicitly quantifies inter-model disagreement. Downstream operators condition their parameter sets and position limits on this distribution rather than on a hard regime label, propagating uncertainty throughout the decision pipeline in a mathematically coherent manner.

The fourth objective—mandatory risk gating—enforces a hard architectural separation between the signal generation layer and the order generation layer. In many production systems this separation is advisory rather than structural: a risk check may be implemented as a function called at the discretion of the strategy author, which may be omitted, misconfigured, or overridden under operational pressure. In this platform, the execution adapter will not accept an order object that does not carry a cryptographically signed approval token issued by the \texttt{RiskOperator} node. The \texttt{RiskOperator} evaluates each proposed action against a hierarchy of constraints: per-instrument notional limits, portfolio-level gross and net exposure bounds, Value-at-Risk and Conditional Value-at-Risk budgets computed under the current regime distribution, maximum allowable drawdown thresholds, and concentration limits expressed as fractions of average daily volume. Approval tokens are time-stamped and expire after a configurable staleness window, ensuring that a signal generated under one risk state cannot be executed after market conditions have materially changed without passing through a fresh risk evaluation cycle.

The fifth objective—provenance tracking—responds to the growing regulatory and operational demand for full explainability of automated trading decisions. Regulators in multiple jurisdictions now require that algorithmic trading firms be able to reconstruct, from durable audit records, the complete causal chain from raw input data to executed order for any transaction under scrutiny. This platform satisfies that requirement by attaching a structured provenance graph to every output produced by every operator node. The provenance graph is a directed acyclic graph whose nodes are immutable data artefacts identified by content-addressed hashes, and whose edges record the operator name, operator version, parameter set hash, and wall-clock execution timestamp of each transformation. This design ensures that any output can be deterministically reproduced from its recorded inputs and operator version, subject only to the availability of the corresponding source data in the data lake. Provenance metadata is stored in a separate append-only ledger, partitioned by trading date, to facilitate efficient retrieval during regulatory examinations without imposing runtime latency on the critical trading path.

The sixth objective—multi-modal operational support—recognises that a quantitative strategy must be validated through a sequence of increasingly demanding testing regimes before it is entrusted with real capital. The platform provides four first-class operational modes that share a common codebase and differ only in their clock and data source implementations. In \emph{historical simulation} mode, a monotonically advancing synthetic clock drives the operator DAG over a fixed historical window, with all data sourced from the immutable data lake; this mode is optimised for throughput and supports vectorised batch evaluation. In \emph{walk-forward backtesting} mode, the historical window is partitioned into an expanding sequence of in-sample training periods and out-of-sample test periods, with model re-fitting occurring at each fold boundary; this mode produces performance statistics that are substantially more predictive of live performance than those produced by a single in-sample backtest. In \emph{paper trading} mode, the system consumes live market data feeds but routes all order objects to a simulated execution engine that applies realistic latency, partial fill, and market impact models rather than forwarding orders to a live broker. In \emph{live trading} mode, orders are forwarded through the Alpaca API adapter with full pre-trade risk checks and post-trade reconciliation. The four modes are distinguished by a single configuration parameter; all operator logic and risk gating code is identical across modes, eliminating the class of bugs that arises when production code diverges from research code.

\section{Document Organisation}

Chapter~2 describes the four-domain financial model, elaborating the conceptual taxonomy that partitions the platform's financial scope into the Earth, Air, Fire, and Water domains, the directed cross-domain influence graph that governs information flow between them, and the formal specification of the cross-domain influence matrix $\Gamma \in \mathbb{R}^{4 \times 4}$ whose entries quantify the directed signal coupling between domain pairs. Chapter~3 presents the operator and agent composition framework in full formal detail, including the typed operator interface, the DAG scheduling semantics, the distinction between stateless and stateful operators, the agent hierarchy that enables multi-level abstraction, and the mechanisms by which operator outputs are versioned and cached. Chapter~4 covers the data pipeline architecture, including the temporal correctness enforcement mechanism, the schema validation layer, the data lake partitioning strategy, the feed normalisation adapters for equity, options, futures, and macroeconomic data sources, and the point-in-time reconstruction service that supports accurate historical simulation.

Chapter~5 details the market regime engine, including the mathematical specification of each constituent classifier, the Bayesian model averaging scheme, the regime vocabulary definition process, the regime-conditional parameter switching mechanism, and the empirical calibration procedures used to ensure that the ensemble posterior is a well-calibrated probability distribution rather than a miscalibrated confidence score. Chapter~6 specifies the risk management framework in full, including the formal constraint hierarchy, the Value-at-Risk and Conditional Value-at-Risk estimation procedures under non-Gaussian return distributions, the cryptographic approval token protocol, the position limit computation methodology, and the drawdown circuit-breaker logic. Chapter~7 covers the execution and simulation layers, including the order routing state machine, the Alpaca API integration, the simulated execution engine's latency and market impact models, and the post-trade reconciliation and attribution pipeline.

Chapter~8 addresses the polyglot implementation architecture, providing a principled justification for the assignment of Python to research and orchestration, C++ to performance-critical numerical computation, Rust to data integrity and storage, and Go to network services and API adapters, together with a description of the language-neutral schema layer that prevents interface divergence across language boundaries. Chapter~9 describes the testing and validation taxonomy, covering unit testing of individual operators, integration testing of operator sub-graphs, statistical validation of regime classifiers, adversarial temporal correctness testing, and the performance benchmarking suite that guards against regression in execution latency. Chapter~10 covers observability, provenance, and audit, including the structured logging schema, the distributed tracing integration, the provenance ledger format, and the regulatory reporting interface. Chapter~11 presents deployment and operational considerations, including the containerisation strategy, the configuration management discipline, the operational runbook for mode transitions, and the incident response procedures for anomalous trading behaviour.

% ============================================================
% CHAPTER 2: FOUR-DOMAIN FINANCIAL MODEL
% ============================================================
\chapter{The Four-Domain Financial Model}

\section{Overview}

The platform organises its financial scope into four conceptual domains: \emph{Earth}, \emph{Air}, \emph{Fire}, and \emph{Water}. These names are not decorative; they encode a deliberate epistemological partition of financial reality into four irreducible analytical stances — long-horizon wealth accumulation, dynamic risk mitigation, volatility-centred alpha extraction, and macro-structural regime awareness. Each domain is a first-class citizen of the system, equipped with its own typed operator algebra, its own agent hierarchy, its own temporal bookkeeping, and its own objective function. The typed operator algebra of each domain is formalised as a monoid $(\mathcal{O}_d, \circ, \mathrm{id}_d)$, where $\mathcal{O}_d$ is the set of well-typed operators for domain $d \in \{\mathrm{Earth}, \mathrm{Air}, \mathrm{Fire}, \mathrm{Water}\}$, $\circ$ denotes operator composition, and $\mathrm{id}_d$ is the identity operator. Composition is subject to a type-compatibility predicate $\tau : \mathcal{O}_d \times \mathcal{O}_d \to \{\top, \bot\}$ that checks input--output type alignment; any composition violating $\tau$ is rejected at parse time by the operator compiler, before any market data is consumed. Each domain's agent hierarchy is stratified into three tiers — \emph{Signal Agents}, \emph{Synthesis Agents}, and \emph{Decision Agents} — corresponding respectively to raw feature extraction, multi-signal aggregation, and actionable portfolio instruction generation. The temporal bookkeeping layer assigns every operator invocation a monotonically increasing logical timestamp drawn from a global Lamport clock, ensuring that causal ordering is preserved even when physical network delays cause out-of-order message delivery between domain subsystems. The objective function of each domain is a domain-specific risk-adjusted return criterion, formally specified in the operator configuration schema and evaluated at every rebalancing epoch by the domain's Decision Agent. Yet the domains do not operate in isolation. They are four projections of a single underlying stochastic process — the evolution of a multi-asset portfolio through time — and the system enforces this unity by making every cross-domain information channel explicit, typed, and subject to the same temporal correctness invariants that govern within-domain computation. No domain is permitted to consume a signal whose timestamp has not been validated by the global clock synchronisation layer, and no domain is permitted to produce an output whose provenance cannot be traced back to a well-formed sequence of operator invocations.

The motivation for a four-domain decomposition, rather than a monolithic strategy engine or an ad hoc collection of signal generators, is both theoretical and operational. Theoretically, the four domains correspond to four distinct risk premia that academic finance has identified as the primary sources of excess return: the equity risk premium (Earth), the insurance premium embedded in options markets (Air), the variance risk premium (Fire), and the macroeconomic cycle premium (Water). Each of these premia has a well-characterised empirical literature. The equity risk premium, documented exhaustively by Mehra and Prescott (1985) and subsequently by the Campbell--Cochrane habit-formation model, represents compensation for bearing undiversifiable aggregate consumption risk over multi-year horizons. The insurance premium embedded in options markets — often called the volatility risk premium in its put-skew manifestation — arises because investors systematically overpay for downside protection relative to the realised frequency of large drawdowns, generating a persistent edge for the disciplined seller of convexity. The variance risk premium, distinct from but related to the insurance premium, is the spread between risk-neutral variance (as measured by the VIX or analogous model-free implied variance estimators) and physical variance (as measured by realised variance over the subsequent period), and has been shown by Carr and Wu (2009) to be a robust, economically significant predictor of short-horizon returns. The macroeconomic cycle premium reflects the time-varying nature of all risk premia as a function of the business cycle, credit cycle, and monetary policy regime, a relationship formalised in the ICAPM tradition by Merton (1973) and extended to dynamic factor models by Chen, Roll, and Ross (1986). Operationally, the decomposition allows specialised agent teams to be assigned to each domain without creating organisational coupling: the volatility team working in the Fire domain need not understand the dividend discount models used by the Earth domain, but the system's cross-domain influence layer ensures that their outputs are correctly weighted and integrated at the portfolio level. This separation of concerns, enforced at the type level and verified by the temporal consistency checker, is the architectural foundation on which the platform's correctness guarantees rest. The type-level enforcement is achieved through a domain-tagged signal type $\mathbf{S}^{(d)}_t$ that carries its domain provenance as a phantom type parameter, making it a compile-time error to pass an Air-domain signal directly to an Earth-domain operator without an explicit cross-domain adapter that is itself registered in the influence graph and subject to audit logging.

\paragraph{Earth Domain.}
The Earth domain embodies the classical long-horizon investor: a rational agent maximising expected utility over a multi-decade horizon under a power utility function $U(W) = W^{1-\gamma}/(1-\gamma)$, where $\gamma > 0$ is the coefficient of relative risk aversion. The domain's primary analytical concern is the construction and maintenance of an equity-centric portfolio whose factor exposures are deliberately calibrated to harvest the equity risk premium and its well-known sub-premia: the value premium (Fama and French, 1992), the profitability premium (Novy-Marx, 2013), the investment premium (Cooper, Gulen, and Schill, 2008), and the low-volatility anomaly (Baker, Bradley, and Wurgler, 2011). The \texttt{PortfolioOpt} operator implements a mean-variance optimisation with Black--Litterman prior views, blending the CAPM equilibrium return vector $\boldsymbol{\Pi} = \delta \boldsymbol{\Sigma} \mathbf{w}_{\mathrm{mkt}}$ (where $\delta$ is the implied risk-aversion coefficient, $\boldsymbol{\Sigma}$ the asset covariance matrix, and $\mathbf{w}_{\mathrm{mkt}}$ the market-capitalisation weights) with proprietary return forecasts $\mathbf{Q}$ via the standard posterior formula $\boldsymbol{\mu}_{\mathrm{BL}} = [(\tau \boldsymbol{\Sigma})^{-1} + \mathbf{P}^\top \boldsymbol{\Omega}^{-1} \mathbf{P}]^{-1}[(\tau \boldsymbol{\Sigma})^{-1} \boldsymbol{\Pi} + \mathbf{P}^\top \boldsymbol{\Omega}^{-1} \mathbf{Q}]$. The \texttt{FactorExposure} operator monitors live factor loadings using a Kalman filter that tracks the time-varying beta coefficients of each holding with respect to the Fama--French--Carhart five-factor model, generating continuous exposure reports consumed by both the Earth domain's own rebalancing scheduler and the Air domain's hedging engine. The \texttt{Rebalancing} operator implements a transaction-cost-aware threshold rebalancing rule: a rebalance is triggered when the $\ell_1$-norm of the deviation of current weights from target weights exceeds a threshold $\epsilon_{\mathrm{reb}}$, which is itself a function of the current bid--ask spread environment as estimated by the \texttt{LiquidityMonitor} sub-operator. Dividend income is tracked separately from capital gain, with reinvestment rules encoded as a domain-level configuration parameter that can be set to immediate reinvestment, scheduled reinvestment, or cash accumulation pending the next rebalancing epoch.

\paragraph{Air Domain.}
The Air domain is the system's risk management nervous system, responsible for the continuous translation of Earth-domain portfolio exposures into a structured programme of derivatives overlays designed to limit drawdown, manage tail risk, and ensure that the total portfolio remains within the risk budget allocated by the Water domain's regime-conditioned capital allocation engine. The Air domain's operators work primarily with listed options and futures, and their correct functioning depends critically on accurate real-time pricing of those instruments. Options pricing within the Air domain uses the Heston (1993) stochastic volatility model as its baseline, with the model parameters $(\kappa, \theta, \sigma_v, \rho, v_0)$ — representing the mean-reversion speed, long-run variance, volatility of volatility, correlation between asset and variance Brownian motions, and initial variance respectively — calibrated daily to the observed implied volatility surface provided by the Fire domain's \texttt{VolSurface} operator. This coupling between the Air and Fire domains is one of the most structurally important linkages in the system: the Air domain consumes the Fire domain's volatility surface in order to price its hedging instruments accurately, while the Fire domain in turn monitors the Air domain's hedging activity to detect order-flow-induced surface distortions that would corrupt its own arbitrage signals. The \texttt{DeltaHedge} operator implements a continuous-time delta-neutral rebalancing strategy under transaction cost constraints, solving the Leland (1985) modified Black--Scholes PDE to determine the optimal rehedging frequency as a function of the bid--ask spread $\xi$ and the option's gamma $\Gamma$: the rehedging interval $\delta t^*$ satisfies $\delta t^* = \xi^2 / (2 \Gamma^2 \sigma^2 S^2)$, balancing the cost of transaction friction against the cost of accumulated delta drift. The \texttt{TailRisk} operator maintains a permanent programme of out-of-the-money put protection, sizing the protective put position using the conditional value-at-risk (CVaR) at the $\alpha = 0.01$ confidence level as the risk measure, and selecting the optimal strike and expiry combination by minimising the Sharpe ratio drag of the protection programme subject to the constraint that the 30-day 99\% CVaR of the total portfolio does not exceed the Water domain's regime-conditioned risk budget $\mathcal{B}^{(r)}_t$.

\paragraph{Fire Domain.}
The Fire domain is the system's volatility intelligence centre, responsible for constructing, maintaining, and exploiting the implied volatility surface of each covered asset. The volatility surface $\sigma_{\mathrm{IV}}(K, T)$ — where $K$ is the option strike and $T$ is the time to expiry — is a rich, high-dimensional object that encodes the market's consensus probability distribution over future asset prices, including all higher moments beyond the variance. The \texttt{VolSurface} operator constructs the surface from quoted option prices using the SVI (Stochastic Volatility Inspired) parameterisation of Gatheral (2004), which represents the total implied variance $w(k) = \sigma_{\mathrm{IV}}^2(k, T) \cdot T$ as a function of log-moneyness $k = \ln(K/F)$ via $w(k) = a + b[\rho(k - m) + \sqrt{(k-m)^2 + s^2}]$, where the parameters $(a, b, \rho, m, s)$ are fitted by least-squares to the observed market quotes subject to the no-calendar-spread and no-butterfly-spread no-arbitrage constraints. The resulting surface is differentiable with respect to both strike and expiry, enabling the \texttt{SkewAnalysis} operator to compute the risk-neutral skewness and excess kurtosis of the return distribution at each expiry, as well as the term structure of the variance risk premium $\mathrm{VRP}(T) = \mathbb{E}^{\mathbb{Q}}[\sigma^2_T] - \mathbb{E}^{\mathbb{P}}[\sigma^2_T]$, where $\mathbb{Q}$ and $\mathbb{P}$ denote the risk-neutral and physical measures respectively. Identifying persistent mispricings in the surface — situations where the model-implied price of a spread or butterfly differs materially from its market quote — is the primary source of alpha for the \texttt{GammaScalp} operator, which constructs delta-neutral, vega-neutral spread positions designed to profit from the convergence of mispriced surface segments toward their arbitrage-free values. All Fire-domain alpha positions are sized using a Kelly-fraction criterion applied to the estimated edge and variance of the convergence trade, with a hard position limit set at $f^* / 2$ (the half-Kelly fraction) to account for estimation error in the edge measurement.

\paragraph{Water Domain.}
The Water domain operates at the longest temporal scale and the broadest information scope of any domain in the system. Its fundamental task is to maintain a probabilistic model of the macroeconomic and market-structural environment — the \emph{regime} — and to communicate that model, in real time, to the parameter-conditioning interfaces of the Earth, Air, and Fire domains. The regime model is formalised as a hidden Markov model (HMM) over a discrete state space $\mathcal{R} = \{r_1, \ldots, r_K\}$, with transition matrix $\mathbf{A} \in [0,1]^{K \times K}$ ($A_{ij} = P(R_{t+1} = r_j \mid R_t = r_i)$) and emission distributions $p(\mathbf{x}_t \mid R_t = r_k)$ parameterised as multivariate Gaussians over a macro feature vector $\mathbf{x}_t$ that includes yield curve slope, credit spreads, equity risk premium estimates, currency volatility, and commodity momentum signals. The \texttt{RegimeDetect} operator implements the Baum--Welch algorithm for offline parameter estimation and the forward algorithm for online filtering, producing a posterior regime distribution $\boldsymbol{\pi}_t = (P(R_t = r_1 \mid \mathbf{x}_{1:t}), \ldots, P(R_t = r_K \mid \mathbf{x}_{1:t}))^\top$ that is updated at every data arrival. The \texttt{MacroFactor} operator decomposes the macro feature vector into orthogonal latent factors using a dynamic factor model estimated by the Kalman smoother, providing a low-dimensional summary of the macroeconomic state that is consumed by the regime model and also broadcast directly to the Earth domain's \texttt{PortfolioOpt} operator as a return-forecasting signal. The \texttt{CrossAsset} operator maintains a dynamic correlation matrix across equity, fixed income, commodity, and currency asset classes, estimated using the DCC-GARCH model of Engle (2002), and uses this matrix to detect cross-asset contagion events — rapid, non-fundamental increases in cross-asset correlation that typically precede or accompany systemic market stress — triggering an immediate risk-off signal to the Air domain's \texttt{TailRisk} operator when the contagion score $\mathcal{C}_t = \|\hat{\boldsymbol{\Sigma}}_t - \hat{\boldsymbol{\Sigma}}_{t-20}\|_F$ exceeds its 95th empirical percentile.

\begin{table}[H]
\centering
\caption{Summary of the four financial domains.}
\label{tab:four_domains}
\small
\begin{tabularx}{\textwidth}{p{2.2cm} p{3.0cm} p{4.5cm} p{4.5cm}}
\toprule
\textbf{Domain} & \textbf{Theme} & \textbf{Primary Focus} & \textbf{Key Operators} \\
\midrule
Earth & Portfolio and Income &
    Equity allocation, dividend income, factor exposures, long-horizon rebalancing &
    \texttt{PortfolioOpt}, \texttt{FactorExposure}, \texttt{Rebalancing} \\
\midrule
Air & Hedging and Risk &
    Options hedging, futures overlays, delta/gamma management, tail risk mitigation &
    \texttt{HedgeOpt}, \texttt{DeltaHedge}, \texttt{TailRisk} \\
\midrule
Fire & Volatility and Alpha &
    Implied volatility surfaces, vol arbitrage, options pricing, tactical alpha &
    \texttt{VolSurface}, \texttt{SkewAnalysis}, \texttt{GammaScalp} \\
\midrule
Water & Macro and Cross-Market &
    Regime classification, macro factor extraction, cross-asset correlation,
    currency and commodity dynamics &
    \texttt{RegimeDetect}, \texttt{MacroFactor}, \texttt{CrossAsset} \\
\bottomrule
\end{tabularx}
\end{table}



\section{Domain Interactions}

The four domains interact through a directed influence graph $\mathcal{G} = (\mathcal{V}, \mathcal{E})$, where the vertex set $\mathcal{V} = \{\mathrm{Earth}, \mathrm{Air}, \mathrm{Fire}, \mathrm{Water}\}$ comprises the four domains and the edge set $\mathcal{E}$ contains a directed edge $(i, j)$ whenever domain $i$ produces a signal that is consumed by domain $j$ as a parameter-conditioning input or a direct data feed. Every edge in $\mathcal{E}$ is annotated with a signal type, a latency bound, and a staleness threshold: if a signal on edge $(i, j)$ has not been refreshed within its staleness threshold, the receiving domain $j$ falls back to a conservative default parameterisation rather than consuming a potentially outdated signal. This fail-safe degradation mechanism ensures that a failure or excessive latency in one domain degrades but does not destroy the correctness of other domains. Let $\hat{y}_t^{(i)}$ denote the aggregate signal produced by domain $i$ at time $t$, where aggregation is performed by that domain's Synthesis Agent using a learned attention-weighted combination of its Signal Agents' outputs. The cross-domain influence matrix $\Gamma \in \mathbb{R}^{4 \times 4}$ is defined entry-wise as:
\begin{equation}
    \Gamma_{ij} = \frac{\partial\, \hat{y}_t^{(i)}}{\partial\, \hat{y}_t^{(j)}}, \qquad i \neq j.
\end{equation}
The diagonal entries $\Gamma_{ii}$ are undefined (a domain does not influence itself through the cross-domain channel, only through its own internal feedback loops), and by convention are set to zero in the stored matrix. Estimation of $\Gamma$ is performed by the \texttt{CrossDomainInfluenceOperator} using rolling-window regression on realised returns, with a window length of 63 trading days (approximately one calendar quarter) to balance responsiveness to structural shifts against estimation noise. The matrix is regime-conditioned, meaning a separate $\Gamma^{(r)}$ is maintained for each market regime $r \in \mathcal{R}$, reflecting the empirical observation that inter-domain dependencies are highly non-stationary: during risk-off regimes, the Water domain's influence on Earth and Air expands dramatically, while during low-volatility, risk-on regimes, the Fire--Air channel carries more information than the Water--Earth channel. The regime-conditioned influence matrix is updated using an exponentially weighted moving average with decay factor $\beta = 0.94$, a standard choice motivated by the RiskMetrics methodology, giving older observations an exponentially decreasing contribution to the current estimate.

\subparagraph{Earth--Air linkage.}
Portfolio allocation decisions in the Earth domain drive hedging requirements in the Air domain through a well-defined demand signal $\mathbf{h}_t^{\mathrm{EA}} \in \mathbb{R}^{N_{\mathrm{hedge}}}$, where each component of $\mathbf{h}_t^{\mathrm{EA}}$ specifies the notional delta exposure in a particular asset that the Air domain must neutralise or partially offset. A large equity overweight — defined as any position in which the Earth domain's equity weight $w_t^{\mathrm{eq}}$ exceeds the Water domain's regime-conditioned equity ceiling $\bar{w}^{(r)}_{\mathrm{eq}}$ by more than the tolerance $\epsilon_w = 0.03$ — triggers a priority request to the \texttt{HedgeOptimizationOperator} for protective puts or collar structures on the equity indices most correlated with the overweight. The \texttt{HedgeOptimizationOperator} solves a constrained optimisation problem that minimises the cost of the hedge (measured in basis points of the notional hedged) subject to the constraint that the post-hedge 30-day CVaR at the 99\% confidence level does not exceed the risk budget $\mathcal{B}^{(r)}_t$ allocated by the Water domain. In addition to this reactive hedging channel, the Earth domain continuously streams its live factor exposures to the Air domain via the \texttt{FactorExposure} operator's broadcast interface, allowing the Air domain to maintain a proactive factor-hedge overlay that is adjusted incrementally at each rebalancing epoch rather than requiring a full hedge reconstruction.

\subparagraph{Water--all-domains linkage.}
The Water domain's regime classifications condition the parameters used by every other domain through a publish--subscribe architecture in which each domain registers as a subscriber to the Water domain's regime event stream. When the posterior regime distribution $\boldsymbol{\pi}_t$ shifts sufficiently — as measured by the Jensen--Shannon divergence $D_{\mathrm{JS}}(\boldsymbol{\pi}_t \| \boldsymbol{\pi}_{t-1}) > \theta_{\mathrm{regime}}$, where $\theta_{\mathrm{regime}}$ is a calibrated sensitivity threshold — the Water domain emits a \texttt{RegimeShiftEvent} carrying the new posterior distribution, the estimated transition probabilities, and the macro factor loadings that drove the shift. Receiving domains respond to this event by reconfiguring their operator parameters: the Earth domain updates its Black--Litterman prior and its factor exposure ceilings; the Air domain updates its hedge ratios, its CVaR risk budget, and its rehedging frequency; the Fire domain updates its volatility surface model priors and its variance risk premium forecasts; and the cross-domain influence matrix $\Gamma^{(r)}$ is rotated to the new regime's estimate. During a high-volatility, risk-off regime — characterised by a VIX above 30, a credit spread widening exceeding 150 basis points, and a negative trailing 60-day equity return — the Earth domain reduces equity allocation toward its regime-conditioned floor $\underline{w}^{(r_{\mathrm{stress}})}_{\mathrm{eq}}$, the Air domain increases hedge ratios by a factor drawn from the regime-specific parameter table, and the Fire domain shifts from net short-volatility strategies (which harvest the variance risk premium in calm markets) to net long-volatility strategies (which benefit from the elevated variance risk premium available when the market is in distress). The Water domain also maintains a cross-asset contagion monitor that operates independently of the regime HMM and provides a faster-acting, model-free early warning signal based on the Frobenius norm of the correlation matrix change $\mathcal{C}_t$, ensuring that sudden market dislocations that have not yet been incorporated into the regime posterior can still trigger a protective response from the Air domain.

\subparagraph{Fire--Air arbitrage.}
Implied volatility mispricings identified by the Fire domain create hedging opportunities in the Air domain through a structured signal relay mediated by the \texttt{VolMispricingBroadcast} interface. When the \texttt{SkewAnalysis} operator detects that the model-implied price of a particular put spread, risk reversal, or butterfly deviates from its market quote by more than a threshold $\delta_{\mathrm{arb}}$ measured in vegas (units of implied volatility), it classifies the mispricing by type — overpriced puts, underpriced calls, excess skew premium, or deficient kurtosis premium — and broadcasts the opportunity to the Air domain. The Air domain's \texttt{HedgeOpt} operator can then exploit this information by constructing hedging structures that simultaneously reduce the portfolio's tail risk and generate a positive expected payoff from the convergence of the mispriced surface segment. Specifically, when the Fire domain identifies overpriced puts on a major equity index — a situation that arises frequently during late-cycle risk-off episodes when retail demand for downside protection drives put implied volatilities above fair value — the Air domain can sell those puts as a funded source of hedge premium, collecting the overpricing while delta-hedging the residual directional risk using futures overlays managed by the \texttt{DeltaHedge} operator. The net effect is a hedge structure whose cost is partially or fully offset by the collected premium, improving the Sharpe ratio drag of the tail risk programme relative to a naive buy-and-hold protective put strategy. The Fire--Air channel is subject to a position-limit constraint enforced at the system level: the notional of put-selling activity initiated through this channel cannot exceed 20\% of the total hedge notional, ensuring that the portfolio does not become systematically short convexity in a way that would defeat the purpose of the Air domain's tail risk programme.

\subparagraph{Earth--Water feedback.}
While the Water domain primarily acts as a parameter-conditioning service for the other domains, the Earth domain provides an important feedback signal that allows the Water domain to calibrate its regime model against observable portfolio outcomes rather than relying exclusively on external macro data. The Earth domain streams its realised factor return attribution to the Water domain at each rebalancing epoch: the vector $\boldsymbol{\phi}_t = (\phi_t^{\mathrm{value}}, \phi_t^{\mathrm{mom}}, \phi_t^{\mathrm{quality}}, \phi_t^{\mathrm{low-vol}}, \phi_t^{\mathrm{size}})^\top$ of factor return contributions provides a real-time cross-section of factor premium realisation that the Water domain's \texttt{MacroFactor} operator uses as an additional emission signal in its HMM, enriching the regime inference with portfolio-level evidence that may lead or lag the macro indicators in the primary feature vector $\mathbf{x}_t$. This feedback loop is carefully designed to avoid look-ahead bias: the factor returns $\boldsymbol{\phi}_t$ are computed from prices at the close of the current period and are only incorporated into the Water domain's posterior at the open of the following period, maintaining the temporal integrity of the entire system.

\section{Cross-Domain Allocation}

The problem of allocating risk capital across the four domains is a portfolio optimisation problem in its own right, distinct from and hierarchically superior to the within-domain allocation problems solved by the individual domain operators. The capital allocated to each domain determines the scale of its operations, the notional of its positions, and ultimately its contribution to the total portfolio's risk-adjusted return. This allocation must be responsive to the current market regime — because different risk premia are more or less accessible and more or less fairly compensated in different regimes — while remaining stable enough that frequent reallocation does not generate excessive transaction costs or destabilise the within-domain strategies that depend on consistent capital availability. The cross-domain allocation problem is therefore formulated as a constrained mean-variance problem in the space of domain weights, using regime-conditioned moments estimated from a combination of historical domain return data and forward-looking structural priors derived from the Water domain's macro factor model.

Let $\mathcal{R} = \{r_1, \ldots, r_K\}$ be the set of $K$ market regimes identified by the Water domain's hidden Markov model, $\boldsymbol{\mu}^{(r)} \in \mathbb{R}^4$ the regime-conditioned expected return vector across the four domains, and $\boldsymbol{\Sigma}^{(r)} \in \mathbb{R}^{4 \times 4}$ the regime-conditioned covariance matrix of domain returns. The expected return vector $\boldsymbol{\mu}^{(r)}$ is estimated by a shrinkage estimator that blends the sample mean of domain returns observed during historical episodes of regime $r$ with a structural prior $\boldsymbol{\mu}_0^{(r)}$ derived from the academic risk premium literature, with shrinkage intensity $\delta^{(r)}$ calibrated to minimise the out-of-sample mean-squared error of the resulting estimator on a held-out validation set. The covariance matrix $\boldsymbol{\Sigma}^{(r)}$ is estimated using the Ledoit--Wolf (2004) analytical shrinkage estimator, which shrinks the sample covariance matrix toward a structured target (the constant-correlation matrix in the regime-conditioned case) by an analytically optimal factor that minimises the Frobenius-norm distance between the estimator and the true covariance matrix under a Gaussian likelihood assumption. The optimal cross-domain allocation weight vector $\boldsymbol{\alpha}^{*}$, constrained to the probability simplex $\Delta^3$ (which enforces non-negativity and full investment), solves:
\begin{equation}
    \boldsymbol{\alpha}^{*} = \arg\max_{\boldsymbol{\alpha} \in \Delta^3}
    \left[
        \boldsymbol{\alpha}^\top \boldsymbol{\mu}^{(r)}
        - \lambda\, \boldsymbol{\alpha}^\top \boldsymbol{\Sigma}^{(r)} \boldsymbol{\alpha}
    \right],
    \label{eq:cross_domain_alloc}
\end{equation}
where $\lambda > 0$ is the risk-aversion parameter, set globally at $\lambda = 2.5$ in the baseline configuration, a value calibrated to produce Sharpe-ratio-maximising allocations under long-run average regime conditions. The quadratic programme in \eqref{eq:cross_domain_alloc} admits a closed-form solution when the simplex constraint is inactive (i.e., when all components of the unconstrained optimum $\boldsymbol{\alpha}^{\mathrm{unc}} = (2\lambda)^{-1}(\boldsymbol{\Sigma}^{(r)})^{-1}\boldsymbol{\mu}^{(r)}$ are non-negative), and is solved via the active-set method otherwise, with warm-starting from the previous period's optimal allocation to accelerate convergence. The solution is updated each time the posterior regime distribution $P(R_t = r_k \mid \mathbf{x}_{1:t})$ shifts materially, as determined by the regime entropy threshold $\hat{U}_t > \theta_U$, where $\hat{U}_t = -\sum_{k=1}^K \pi_{t,k} \ln \pi_{t,k}$ is the Shannon entropy of the posterior at time $t$ and $\theta_U$ is set to the 80th percentile of the historical entropy distribution. Between regime-triggered recomputations, the allocation vector $\boldsymbol{\alpha}_t$ drifts passively with the domain returns and is subject to a tolerance-band rebalancing rule: if $\|\boldsymbol{\alpha}_t - \boldsymbol{\alpha}^*\|_1 > \epsilon_{\boldsymbol{\alpha}}$, where $\epsilon_{\boldsymbol{\alpha}} = 0.05$ is the rebalancing tolerance, a rebalancing trade is executed to restore $\boldsymbol{\alpha}_t$ to $\boldsymbol{\alpha}^*$. This combination of regime-triggered reoptimisation and tolerance-band drift management ensures that the cross-domain allocation remains strategically aligned with the current market environment while keeping rebalancing costs at an operationally acceptable level.

% ============================================================
% CHAPTER 3: OPERATOR AND AGENT COMPOSITION FRAMEWORK
% ============================================================
\chapter{Operator and Agent Composition Framework}

\section{Core Abstractions}

The platform's execution model is built on three nested abstractions:
Operators, Agents, and Systems.

\begin{definition}[Operator]
An \emph{Operator} is a typed, stateless, referentially transparent function:
\begin{equation}
    O : \mathcal{S} \times \mathcal{I}_t \rightarrow \mathcal{S}',
\end{equation}
where $\mathcal{S}$ is the shared state space, $\mathcal{I}_t$ is the
information set legitimately available at time $t$, and $\mathcal{S}'$ is
the output state space. An Operator must satisfy:
\begin{enumerate}
    \item \emph{Temporal validity}: $O(s, \mathcal{I}_t) \in \sigma(\mathcal{I}_t)$
          for all $s \in \mathcal{S}$, $t \in \mathcal{T}$.
    \item \emph{Referential transparency}: identical inputs produce identical outputs.
    \item \emph{Failure transparency}: errors are surfaced as typed failure objects,
          not silently swallowed.
\end{enumerate}
\end{definition}

\begin{definition}[Agent]
An \emph{Agent} is a named, versioned composition of Operators:
\begin{equation}
    A = (O_1, O_2, \ldots, O_k, \mathbf{w}),
\end{equation}
where $\mathbf{w} \in \Delta^{k-1}$ is a weight vector on the $(k-1)$-dimensional
probability simplex. The aggregated output is:
\begin{equation}
    \hat{y}_t = \sum_{i=1}^{k} w_i \cdot O_i(s, \mathcal{I}_t).
\end{equation}
\end{definition}

\begin{definition}[System]
A \emph{System} is an orchestrated collection of Agents whose dependency
graph forms a directed acyclic graph (DAG). A topological execution order
$\tau : V \rightarrow \{1, \ldots, |V|\}$ satisfies:
\begin{equation}
    (O_i, O_j) \in E \implies \tau(O_i) < \tau(O_j).
\end{equation}
\end{definition}

\section{Operator Taxonomy}

\begin{table}[H]
\centering
\caption{Primary operator categories and their responsibilities.}
\label{tab:operator_taxonomy}
\small
\begin{tabularx}{\textwidth}{p{4.25cm} p{10.5cm}}
\toprule
\textbf{Operator Category} & \textbf{Responsibility} \\
\midrule
\texttt{DataOperator} &
    Ingests and validates raw market data from exchanges, data vendors,
    and alternative data sources. Enforces schema compliance and
    provenance tagging. \\
\midrule
\texttt{FeatureOperator} &
    Transforms validated, normalised data into typed feature vectors.
    Includes technical indicators, candlestick pattern detectors,
    volatility estimators, and factor exposures. \\
\midrule
\texttt{MarketRegimeOperator} &
    Classifies the current market state into a probability distribution
    over regime labels. Outputs \texttt{MarketRegimeState} with
    posterior probabilities, confidence, and uncertainty. \\
\midrule
\texttt{SignalOperator} &
    Generates directional or positioning signals in the range $[-1, 1]$
    from feature vectors and regime state. Each signal carries an
    associated confidence score. \\
\midrule
\texttt{RiskOperator} &
    Evaluates proposed \texttt{OrderIntent} objects against multi-dimensional
    risk constraints. May reject, resize, delay, or hedge any intent. \\
\midrule
\texttt{PortfolioOperator} &
    Translates approved signals and risk assessments into target portfolio
    weights. Enforces diversification, leverage, and concentration constraints. \\
\midrule
\texttt{ExecutionOperator} &
    Simulates or routes order execution. Initial implementation targets
    paper execution only. \\
\midrule
\texttt{ValidationOperator} &
    Enforces temporal correctness, schema compliance, and invariant
    preservation at each pipeline stage. \\
\midrule
\texttt{NewsOperator} &
    Ingests, classifies, and extracts quantitative signals from
    unstructured news and earnings communications. \\
\midrule
\texttt{SimulationOperator} &
    Manages historical replay, Monte Carlo scenario generation,
    and walk-forward test orchestration. \\
\bottomrule
\end{tabularx}
\end{table}

\section{Operator Composition Algebra}

Sequential composition of two operators $O_1$ and $O_2$ is defined as:
\begin{equation}
    (O_1 \circ O_2)(s, \mathcal{I}_t) = O_1\!\left(O_2(s, \mathcal{I}_t),\, \mathcal{I}_t\right).
\end{equation}
Composition is associative:
\begin{equation}
    (O_1 \circ O_2) \circ O_3 = O_1 \circ (O_2 \circ O_3).
\end{equation}
A valid pipeline $\Pi = O_n \circ \cdots \circ O_1$ satisfies the
temporal validity invariant:
\begin{equation}
    \Pi(s, \mathcal{I}_t) \in \sigma(\mathcal{I}_t) \quad \forall\, s \in \mathcal{S},\; t \in \mathcal{T}.
\end{equation}

\section{The Full Operator Pipeline}

The complete operator pipeline from raw data to risk-gated execution
decision is illustrated conceptually in Table~\ref{tab:pipeline_stages}.

\begin{table}[H]
\centering
\caption{End-to-end operator pipeline stages.}
\label{tab:pipeline_stages}
\small
\begin{tabularx}{\textwidth}{p{1.5cm} p{3.8cm} p{10.0cm}}
\toprule
\textbf{Stage} & \textbf{Operator} & \textbf{Output Type and Description} \\
\midrule
1 & \texttt{DataOperator} &
    \texttt{ValidatedMarketData} --- schema-validated, provenance-tagged
    raw observations. \\
\midrule
2 & \texttt{NormalisationOperator} &
    \texttt{NormalisedData} --- adjusted prices, aligned timestamps,
    currency-normalised values, missing-value-handled series. \\
\midrule
3 & \texttt{FeatureOperator} &
    \texttt{FeatureVector} --- technical indicators, volatility estimates,
    pattern presence flags, factor exposures. \\
\midrule
4 & \texttt{MarketRegimeOperator} &
    \texttt{MarketRegimeState} --- posterior probability vector over
    regime labels, regime entropy $H_t$, normalised uncertainty $\hat{U}_t$. \\
\midrule
5 & \texttt{SignalOperator} &
    \texttt{Signal} --- directional value in $[-1, 1]$ with confidence
    $c \in (0, 1]$ and provenance metadata. \\
\midrule
6 & \texttt{PortfolioOperator} &
    \texttt{OrderIntent} --- proposed target weights or order sizes
    arising from the signal and current portfolio state. \\
\midrule
7 & \texttt{RiskOperator} &
    \texttt{RiskDecision} --- approved, rejected, resized, or hedged
    \texttt{OrderIntent} with detailed risk metadata. \\
\midrule
8 & \texttt{ExecutionOperator} &
    \texttt{ExecutionRecord} --- simulated or live fill with timestamps,
    costs, and a correlation ID linking back to stage~1. \\
\midrule
9 & \texttt{ValidationOperator} &
    Cross-stage invariant check; emits \texttt{AuditRecord} for every
    stage transition. \\
\bottomrule
\end{tabularx}
\end{table}

% ============================================================
% CHAPTER 4: DATA PIPELINE ARCHITECTURE
% ============================================================
\chapter{Data Pipeline Architecture}

\section{Data Categories}

The platform ingests, normalises, and version-controls data across six primary categories, each of which imposes distinct contractual requirements on latency, quality assurance, and temporal alignment. The heterogeneity of these categories is a deliberate design feature rather than an accidental consequence of market structure: different asset classes trade on different venues, publish data under different licensing regimes, and update at frequencies spanning several orders of magnitude. A unified ingestion framework must therefore treat each category as a first-class citizen with its own \texttt{IngestionPolicy}, specifying acceptable staleness bounds, quality-flag thresholds, and the identity of the upstream \texttt{NormalisationOperator} responsible for transforming raw vendor payloads into the canonical internal schema.

\textbf{Price and volume data} constitutes the highest-frequency and lowest-latency category ingested by the platform. Raw market data is consumed at the tick level from direct exchange feeds and consolidated tape providers, capturing every trade print and every change to the national best bid and offer (NBBO). From these atomic events the platform derives a hierarchy of OHLCV bar aggregates at standardised frequencies: 1-second, 1-minute, 5-minute, 15-minute, 1-hour, and daily. Each bar is stamped with both the exchange timestamp and the local ingestion timestamp, and the two are reconciled to detect clock skew in excess of configurable tolerances. Order-book snapshots at configurable depth levels (L1, L2, and L3 where available) are maintained in a rolling buffer, enabling microstructure features such as bid-ask spread, depth imbalance, and weighted mid-price to be computed on demand by downstream \texttt{FeatureOperator} instances. All price data is adjusted for corporate actions --- splits, dividends, rights issues --- using a deterministic adjustment chain that is itself versioned, so that any historical price series can be reproduced exactly from its adjustment-chain version identifier.

\textbf{Derivatives data} introduces additional dimensionality relative to equities because each underlying asset is associated with a potentially large and dynamically evolving set of listed instruments. Options chains are ingested at the full strike-expiry grid, with implied volatility surfaces constructed by the \texttt{VolatilitySurfaceOperator} using a no-arbitrage interpolation scheme that enforces calendar-spread monotonicity and butterfly positivity. The resulting surface is parameterised as a SSVI (Surface Stochastic Volatility Inspired) functional form, enabling efficient interpolation and extrapolation while respecting the static arbitrage constraints derived by Gatheral and Jacquier. Futures curves are handled separately, with explicit modelling of roll dates, basis, and carry. Greeks (delta, gamma, vega, theta, rho, and higher-order cross-derivatives) are computed analytically where closed-form expressions exist and via automatic differentiation for more complex payoff structures. All derivatives data objects carry a reference to the underlying asset identifier, the pricing model version, and the interest-rate curve used, ensuring full reproducibility of any computed quantity.

\textbf{Macroeconomic data} differs fundamentally from market data in that its release schedule is known in advance but its values are not, and in that its impact on asset prices is typically non-linear and regime-dependent. The platform maintains a structured economic calendar that records, for each scheduled release, the consensus forecast, the previous realised value, the revision history, and the release timestamp down to the millisecond. Surprise factors --- the standardised deviation of the actual release from the consensus --- are computed immediately upon ingestion and propagated as derived features to all dependent operators. Central bank communications, including meeting minutes, speeches, and policy statements, are ingested as raw text and processed by the \texttt{NLPOperator} to extract sentiment scores, topic vectors, and forward-guidance signals. Yield curve data is represented as a Nelson-Siegel-Svensson parameterisation updated at each tenor point revision, with the full parameter vector and its covariance structure stored alongside the raw zero-coupon rates.

\textbf{Fundamental data} is characterised by its low update frequency, its high information content, and the significant complexity of its accounting interpretation. Financial statements --- income statements, balance sheets, and cash flow statements --- are ingested from both primary sources (EDGAR filings, exchange announcements) and normalised vendor feeds, with systematic reconciliation between the two to detect vendor errors or restatements. The platform applies a configurable set of accounting normalisation adjustments, including operating-lease capitalisation, R\&D capitalisation, and pension liability recognition, each of which is independently versioned. Point-in-time databases are maintained to ensure that backtests only use information that was publicly available at the time of the simulated decision: restated figures are stored with their restatement dates, and the original as-reported values are preserved in a separate column family. Analyst estimate revisions, earnings surprise histories, and guidance changes are tracked at the individual analyst level where data permits, enabling the construction of estimate dispersion and revision momentum features.

\textbf{Alternative data} encompasses the rapidly expanding universe of non-traditional information sources that carry predictive signal orthogonal to conventional financial data. Satellite imagery of retail car parks, industrial facilities, and agricultural land is processed by computer vision pipelines to generate foot-traffic counts, capacity utilisation estimates, and crop yield forecasts. Anonymised and aggregated credit card transaction data is ingested at weekly intervals and normalised to same-store revenue proxies on a per-merchant and per-sector basis. Web traffic data from panel providers is used to construct digital engagement metrics for e-commerce and SaaS companies. Social sentiment is derived from large-scale analysis of financial discussion forums, news comment sections, and professional networks, with entity-linking to map free-text mentions to canonical security identifiers. Supply chain indicators, including shipping container tracking data and procurement tender announcements, are processed to construct network-graph representations of inter-company dependencies. Because alternative data carries substantially higher legal and ethical risk than exchange data, every alternative data source is subject to a \texttt{DataGovernanceReview} before ingestion, and its license terms are stored as a structured metadata object alongside every derived feature.

\textbf{News and events data} bridges the structured world of corporate disclosures and the unstructured world of real-time narrative. Earnings announcements and earnings call transcripts are ingested immediately upon release and processed through a multi-stage NLP pipeline that extracts quantitative guidance revisions, management sentiment, and question-and-answer sentiment asymmetry. Regulatory filings, including 8-Ks, 13-Fs, 13-Ds, and SEC comment letters, are parsed with schema-aware extractors that map filing sections to structured data fields. Geopolitical event streams are sourced from specialised providers and classified by event type, geographic scope, and estimated market impact severity. Technology news is monitored for product launches, patent grants, partnership announcements, and executive changes, with each event linked to one or more affected security identifiers. All events are stored in a unified event store with a common schema that records the event type, the affected entities, the publication timestamp, the ingestion timestamp, and a confidence score reflecting the extraction quality.



\section{Provenance Requirements}

Every datum used by any decision operator must carry the following provenance
attributes:

\begin{table}[H]
\centering
\caption{Required provenance attributes for all data objects.}
\label{tab:provenance}
\small
\begin{tabularx}{\textwidth}{p{4.75cm} X}
\toprule
\textbf{Attribute} & \textbf{Description} \\
\midrule
\texttt{source} & Originating data vendor or exchange identifier. \\
\texttt{source\_timestamp} & The timestamp assigned by the data source. \\
\texttt{ingestion\_timestamp} & The UTC wall-clock time at which the datum
    was ingested by the platform. \\
\texttt{normalisation\_version} & The version string of the normalisation
    operator that processed the datum. \\
\texttt{transformation\_version} & The version string of any feature-engineering
    transformation applied downstream. \\
\texttt{quality\_status} & Enumerated quality flag: \texttt{CLEAN},
    \texttt{INTERPOLATED}, \texttt{STALE}, \texttt{SUSPECT}, \texttt{REJECTED}. \\
\texttt{correlation\_id} & A UUID linking this datum to the decision record
    it ultimately contributed to. \\
\bottomrule
\end{tabularx}
\end{table}

\section{Temporal Alignment}

Financial data sources operate on incompatible temporal grids. Equity
prices are quoted on exchange calendars; futures roll on specific dates;
macroeconomic releases arrive at irregular intervals; alternative data
has its own update cadence. The \texttt{TemporalAlignmentOperator} is
responsible for:


\paragraph{UTC Canonicalisation and Time Zone Provenance.}
The first and most foundational responsibility of the \texttt{TemporalAlignmentOperator} is the
conversion of every ingested timestamp into Coordinated Universal Time (UTC), accompanied by the
explicit annotation of the originating time zone and a boolean flag indicating whether Daylight
Saving Time (DST) was in effect at the moment of observation. This requirement is non-trivial in
practice: financial data vendors routinely publish timestamps in local exchange time, which
introduces a two-dimensional ambiguity — the offset from UTC is both jurisdiction-dependent and
seasonally variable. For example, the New York Stock Exchange operates on Eastern Standard Time
(UTC$-$5) during winter months and Eastern Daylight Time (UTC$-$4) during summer, while the
London Stock Exchange alternates between GMT (UTC$+$0) and BST (UTC$+$1). Failing to account for
these transitions produces systematic one-hour displacements in aligned time series that are
notoriously difficult to detect post-hoc. The operator therefore maintains a versioned, auditable
time zone database (conforming to the IANA \texttt{tzdata} standard) and applies it deterministically
at ingestion time. The resulting metadata tuple $(\tau_{\mathrm{UTC}},\, \mathrm{tz\_origin},\,
\mathrm{dst\_flag})$ is persisted alongside the datum and is immutable for the lifetime of the
record. Any downstream reprocessing that alters the canonical UTC value must produce a new record
version rather than overwriting the original, thereby preserving the complete audit trail.

\paragraph{Frequency Grid Alignment via Causal Forward-Fill.}
Having canonicalised all timestamps to UTC, the operator must reconcile the heterogeneous update
cadences of disparate data sources onto a single, uniform target frequency grid $\mathcal{G}_\Delta$,
where $\Delta$ denotes the grid resolution (e.g., $\Delta = 1\,\text{minute}$, $\Delta =
1\,\text{day}$). The alignment procedure follows strictly causal, forward-fill semantics: for any
grid point $t_k \in \mathcal{G}_\Delta$, the value assigned to a series $X$ is the most recently
observed value $X_{\tau}$ such that $\tau \leq t_k$. Formally, the aligned value is:
\begin{equation}
    \tilde{X}_{t_k} \;=\; X_{\tau^*}, \quad \text{where} \quad
    \tau^* = \sup\bigl\{\tau \in \mathcal{T}_X : \tau \leq t_k\bigr\},
\end{equation}
where $\mathcal{T}_X$ is the set of timestamps at which source $X$ was actually observed. This
construction guarantees that $\tilde{X}_{t_k}$ is measurable with respect to $\mathcal{I}_{t_k}$
and introduces no look-ahead bias. Any interpolation scheme that uses future observations — linear
interpolation being the most prevalent offender — is categorically prohibited because it implicitly
conditions on $\mathcal{I}_{t_{k+j}}$ for $j \geq 1$, constituting a temporal leakage violation.
The operator is further responsible for handling irregularly spaced source grids, such as
macroeconomic release calendars or corporate earnings announcement schedules, by treating them as
special cases of the general forward-fill mechanism with potentially large inter-observation gaps.
The filled values carry a provenance annotation indicating that they derive from a prior observation
rather than a contemporaneous one, enabling downstream consumers to condition their behaviour on
data freshness when necessary.

\paragraph{Staleness Detection and Quality Flag Propagation.}
Not all silences in a time series are semantically equivalent. A forward-filled value that is one
tick old in a high-frequency equity series is qualitatively different from a macroeconomic
indicator that has not been refreshed in several days. The \texttt{TemporalAlignmentOperator}
therefore implements a configurable staleness threshold $\delta_{\max}$, specified per data
source and per asset class, such that any aligned grid point $t_k$ for which
$t_k - \tau^* > \delta_{\max}$ receives a \texttt{STALE} quality flag rather than a
\texttt{CLEAN} one. More precisely, let $\Delta\tau_k = t_k - \tau^*$ denote the age of the
most recent observation at grid point $t_k$. Then:
\begin{equation}
    \mathrm{quality\_status}_{t_k} =
    \begin{cases}
        \texttt{CLEAN}  & \text{if } \Delta\tau_k \leq \delta_{\max}, \\
        \texttt{STALE}  & \text{if } \Delta\tau_k > \delta_{\max}.
    \end{cases}
\end{equation}
The threshold $\delta_{\max}$ is not a single global constant but a source-specific configuration
parameter stored in the operator's versioned parameter manifest. For liquid equity mid-prices,
a reasonable default might be $\delta_{\max} = 5\,\text{minutes}$; for daily macroeconomic
series, $\delta_{\max}$ might be set to $3\,\text{business days}$. Critically, the \texttt{STALE}
flag is propagated transitively: if a derived feature $F_t = f(\tilde{X}_t, \tilde{Y}_t)$ depends
on any input carrying a \texttt{STALE} flag, the feature itself inherits the worst-case quality
status of its inputs. This monotone degradation of quality metadata ensures that no derived signal
can silently present stale constituent data as fresh, and that portfolio construction and risk
management layers can make informed decisions about whether to trust any given feature value.

\paragraph{Strict Enforcement of the Information Boundary.}
The final and most rigorous responsibility of the \texttt{TemporalAlignmentOperator} is the
enforcement of the hard causal information boundary: no datum whose source timestamp $\tau_s$ is
strictly greater than the current processing time $t$ may appear in the information set
$\mathcal{I}_t$ made available to any downstream operator at time $t$. Formally, this is the
constraint:
\begin{equation}
    \forall\, d \in \mathcal{I}_t :\quad \tau_s(d) \leq t,
\end{equation}
where $\tau_s(d)$ denotes the source-assigned timestamp of datum $d$. This constraint must be
enforced both in live trading and — critically — during historical backtesting, where the
temptation to use the full dataset as a convenient lookup table is a persistent source of
inadvertent leakage. The operator implements this boundary through a \emph{point-in-time cursor}
mechanism: for any query at logical time $t$, the operator materialises the state of every source
series as it would have appeared at time $t$, using only observations with $\tau_s \leq t$.
This is technically equivalent to maintaining a bi-temporal data model in which one temporal axis
records the source timestamp and the other records the ingestion timestamp, and queries are always
evaluated against a snapshot of the former axis at or before $t$. Violations of this constraint
are not merely statistical errors; they constitute a fundamental breach of the strategy's causal
validity and will produce backtested performance figures that are irreproducible in live deployment.
The platform enforces this invariant at the infrastructure level rather than relying on operator
discipline, making it structurally impossible for any consumer of $\mathcal{I}_t$ to access a
datum with $\tau_s > t$ regardless of the query pattern employed.

\section{The Temporal Leakage Invariant}

Let $\mathcal{I}_t$ denote the information set legitimately available at
time $t$. A strategy $\pi$ is \emph{temporally valid} if and only if:
\begin{equation}
    \forall\, t :\quad \text{Decision}_t(\pi) \in \sigma(\mathcal{I}_t).
\end{equation}
A \emph{temporal leakage violation} occurs when:
\begin{equation}
    \exists\, t,\, \tau > 0 :\quad \text{Decision}_t(\pi)
    \text{ depends on } \mathcal{I}_{t+\tau} \setminus \mathcal{I}_t.
\end{equation}

The most common sources of leakage are documented in Table~\ref{tab:leakage_sources}.

\begin{table}[H]
\centering
\caption{Common sources of temporal data leakage and mitigations.}
\label{tab:leakage_sources}
\small
\begin{tabularx}{\textwidth}{p{4.5cm} p{9.5cm}}
\toprule
\textbf{Leakage Source} & \textbf{Mitigation} \\
\midrule
Future price levels used in normalisation &
    Normalise using only the rolling statistics available up to $t-1$. \\
\midrule
Revised macroeconomic data &
    Use point-in-time (vintage) data series, not the latest revised release. \\
\midrule
Survivorship bias &
    Include all instruments that were listed at time $t$, including
    subsequently delisted securities. \\
\midrule
Train/validation contamination &
    Enforce a purging window and embargo period around all walk-forward splits. \\
\midrule
Improperly adjusted prices &
    Apply adjustments (splits, dividends) using only information known at $t$. \\
\midrule
Future news leakage &
    Gate news data by ingestion timestamp, not publication date of the
    post-processed version. \\
\bottomrule
\end{tabularx}
\end{table}

\section{Formal Proof of Temporal Validity and Leakage Consequences}

The temporal leakage invariant stated above is not merely a software engineering convention — it carries profound implications for the statistical validity of any empirical performance evaluation and for the theoretical correctness of the strategy as a decision-making process. To make this precise, consider the filtered probability space $(\Omega, \mathcal{F}, \{\mathcal{F}_t\}_{t \geq 0}, \mathbb{P})$, where the filtration $\{\mathcal{F}_t\}$ represents the natural information flow of the market. A trading strategy $\pi$ is said to be \emph{adapted} to this filtration if and only if every decision $\text{Decision}_t(\pi)$ is measurable with respect to $\mathcal{F}_t$, meaning it depends only on information that has been revealed by time $t$. The condition $\text{Decision}_t(\pi) \in \sigma(\mathcal{I}_t)$ is the discretised, operational counterpart of this measurability requirement, where $\sigma(\mathcal{I}_t)$ denotes the $\sigma$-algebra generated by the observation set $\mathcal{I}_t$. Violating this condition by allowing $\text{Decision}_t(\pi)$ to depend on elements of $\mathcal{I}_{t+\tau} \setminus \mathcal{I}_t$ for any $\tau > 0$ renders the strategy non-adapted, which immediately invalidates the martingale and semi-martingale representations on which the majority of quantitative finance theory rests. In particular, the fundamental theorem of asset pricing, the no-arbitrage conditions expressed via equivalent martingale measures, and the dynamic programming decompositions underlying optimal control formulations of portfolio selection all presuppose adapted strategies; a non-adapted strategy can manufacture arbitrage in a frictionless, complete market purely through its access to future information, producing P\&L that is economically meaningless and statistically non-replicable.

The practical consequences of leakage are equally severe and are often more insidious than the theoretical concerns because they manifest as inflated but plausible-looking backtested returns rather than as obviously impossible results. Consider a normalisation scheme in which each raw price series $p_t^{(i)}$ is scaled by its full-sample standard deviation $\hat{\sigma}^{(i)} = \left(\frac{1}{T}\sum_{t=1}^T (p_t^{(i)} - \bar{p}^{(i)})^2\right)^{1/2}$, where $T$ is the length of the entire historical dataset. This quantity is computed using observations from all $T$ time steps, including those at $t' > t$ for any intermediate evaluation point $t < T$. A signal derived from this normalised series at time $t$ therefore implicitly encodes information about the future volatility of the series beyond $t$, which constitutes a leakage of type one in Table~\ref{tab:leakage_sources}. The correct procedure is to replace $\hat{\sigma}^{(i)}$ with the rolling estimator $\hat{\sigma}_t^{(i)} = \left(\frac{1}{w}\sum_{s=t-w}^{t-1} (p_s^{(i)} - \bar{p}_{t-1}^{(i)})^2\right)^{1/2}$ for some lookback window $w$, which depends only on observations available in $\mathcal{I}_t$. The difference in signal quality between these two schemes is not merely a numerical nuance: in periods of structural volatility shifts — such as the transition from the low-volatility regime of 2012–2019 to the elevated volatility of 2020 — the full-sample estimator retrospectively normalises the pre-transition data by a quantity that was unknowable before the transition occurred, systematically biasing signal magnitudes and the apparent Sharpe ratio of any strategy that consumes them.

The point-in-time cursor mechanism implemented by the \texttt{TemporalAlignmentOperator} provides the infrastructure-level guarantee that prevents all of these failure modes simultaneously. The key design insight is that temporal validity cannot be delegated to operator discipline or code review alone; it must be enforced by a mechanism that makes violations structurally impossible rather than merely inadvisable. The bi-temporal data model underlying the cursor maintains two independent time axes for every stored observation: the \emph{event time} $\tau_e$, which records when the underlying event occurred in the market, and the \emph{ingestion time} $\tau_i$, which records when the datum arrived in the platform's data store. A point-in-time query at logical time $t$ retrieves only observations for which $\tau_e \leq t$ \emph{and} $\tau_i \leq t$, ensuring that even data describing past events is excluded if it was not yet available to the platform at time $t$. This is particularly critical for macroeconomic releases, which are subject to multi-vintage revision cycles spanning months or years: the GDP growth rate for Q3 of year $Y$ may be initially reported in late Q4 of year $Y$, revised in Q1 of year $Y+1$, and substantially revised again in Q2 of year $Y+2$. A naive query returning the ``current'' value for Q3 of year $Y$ would retrieve the final revised figure, which was unavailable at any point before Q2 of year $Y+2$, introducing a leakage window of up to two years. The bi-temporal model eliminates this class of violation by construction.

The formal leakage audit procedure executed at strategy registration time works as follows. For every operator $\phi_i$ in the computation graph $\mathcal{G} = (\mathcal{V}, \mathcal{E})$, the platform enumerates the complete set of data source identifiers in $\phi_i$'s declared dependency manifest and verifies, for each source $s_j$, that the timestamp function $\tau_s : s_j \to \mathbb{R}$ satisfies $\tau_s(d) \leq t$ for all observations $d$ retrieved at processing time $t$. This verification is performed symbolically at registration time through a static analysis of the operator's data access patterns, and is additionally enforced dynamically at runtime through interceptors placed on all data access methods. Any operator that requests a datum with $\tau_s > t$ raises a \texttt{TemporalIntegrityViolation} exception, which is propagated as an unrecoverable fault — the strategy is halted and the violation is logged with full provenance information for compliance review. This zero-tolerance policy ensures that the temporal validity guarantee is maintained not only during initial development and backtesting, but also throughout the entire operational lifecycle of the strategy as it is modified, extended, and subjected to regime changes in the underlying data infrastructure.

The distinction between survivorship bias and point-in-time data integrity deserves particular elaboration because the two failure modes, while structurally similar, arise through different mechanisms and require different mitigations. Survivorship bias occurs when the universe of instruments evaluated at time $t$ during backtesting consists only of instruments that were still listed at the time the backtest was constructed, retroactively excluding all instruments that were delisted, went bankrupt, were acquired, or were otherwise removed from the market between time $t$ and the present. The consequence is a systematic upward bias in the apparent quality of the investment universe, because the excluded instruments are disproportionately those that experienced adverse outcomes. If the strategy has any mechanism for selecting among instruments — even a passive market-cap weighting — the exclusion of historically poor performers biases the selection toward the survivors, producing an apparent alpha that is actually a statistical artefact of the evaluation methodology. The platform's instrument universe management module addresses this through a \emph{constituent snapshot} data structure: at each time $t$, the platform reconstructs the exact set of instruments that satisfied the listing criteria at time $t$, including all subsequently delisted securities, using the bi-temporal ingestion records that capture when each instrument entered and exited each index or eligibility screen. The formal requirement is that the universe $\mathcal{U}_t$ satisfies $\mathcal{U}_t = \{i : \tau_{\text{list}}^{(i)} \leq t \leq \tau_{\text{delist}}^{(i)}\}$, where $\tau_{\text{list}}^{(i)}$ and $\tau_{\text{delist}}^{(i)}$ are the listing and delisting timestamps of instrument $i$, respectively, both recorded at their original event times rather than at any subsequent revision or discovery time.

The train/validation contamination leakage listed in Table~\ref{tab:leakage_sources} is perhaps the most technically subtle entry, because it arises not from direct access to future data but from the indirect leakage of future information through the model selection and hyperparameter tuning process. In a standard walk-forward validation framework, the dataset is partitioned into a sequence of training windows $\mathcal{T}_1, \mathcal{T}_2, \ldots, \mathcal{T}_n$ and corresponding validation windows $\mathcal{V}_1, \mathcal{V}_2, \ldots, \mathcal{V}_n$, with each validation window immediately following its corresponding training window. The purging requirement stipulates that no observation within a purging window of half-width $\delta_p$ around the boundary between $\mathcal{T}_k$ and $\mathcal{V}_k$ may be used in either the training or the validation set for fold $k$. The theoretical justification is that many financial time series exhibit label overlap: if the target variable $y_t$ for a return prediction model is the $h$-day ahead return $\sum_{s=1}^{h} r_{t+s}$, then the labels $y_t$ and $y_{t+1}$ share $h-1$ overlapping return observations. Including $y_{t+1}$ in the training set while evaluating on $y_t$ during validation introduces an information channel through the shared return observations, which inflates the apparent out-of-sample performance of the model by an amount that grows with the overlap fraction $1 - 1/h$. The embargo period $\delta_e$ further extends the exclusion zone beyond the training window's end, preventing the model from capturing short-range autocorrelation in the residuals as if it were genuine predictive signal. The combined purge-embargo requirement mandates that validation fold $k$ exclude from training all observations within the interval $[\min(\mathcal{V}_k) - \delta_p - \delta_e,\, \min(\mathcal{V}_k) + \delta_p]$, where the values of $\delta_p$ and $\delta_e$ are set to at least $h$ and $2h$ bars respectively for a prediction horizon of $h$ bars, as derived from the overlap analysis of \citet{deMLforAF2018}.



% ============================================================
% CHAPTER 5: MARKET REGIME ENGINE
% ============================================================
\chapter{Market Regime Engine}

\section{Motivation}

Market regime is the single most important conditioning variable in any quantitative strategy, yet it is also among the most frequently neglected. The failure mode is insidious: a strategy calibrated exclusively in one market environment accumulates risk exposures that are entirely invisible until the environment shifts, at which point the strategy may suffer rapid, severe drawdowns that are structurally impossible to hedge within the strategy's own parameter space. A momentum strategy that performs well in trending, low-volatility markets becomes a reliable loss generator in choppy, mean-reverting environments because the directional persistence assumption embedded in its signal construction is simply no longer satisfied by the data-generating process. A short-volatility strategy that earns steady option premium in calm periods suffers catastrophic, often gap-risk-induced losses during volatility spikes precisely because the return distribution's tail mass expands far beyond the Gaussian approximation on which most option pricing models are implicitly conditioned. The cross-sectional dispersion of individual equity returns, the autocorrelation structure of price increments, the co-integration relationships between related assets, and even the microstructure properties of bid-ask spreads all shift materially across regimes. Without an explicit, probabilistically coherent representation of the current regime, strategy parameters optimised in one regime are silently misapplied in another, producing what practitioners sometimes call ``regime contamination'' of the signal-to-noise ratio.

The platform therefore elevates regime identification to a first-class architectural concern, treating it not as an optional enrichment layer but as the foundational conditioning variable through which every downstream operator---signal generation, portfolio optimisation, execution sizing, and risk limit selection---is filtered. This design choice reflects a fundamental epistemic commitment: uncertainty about which market regime is currently active must be represented explicitly as a probability distribution, not collapsed into a single point estimate, because the cost of a regime misclassification is asymmetric and non-linear. Propagating the full posterior $\boldsymbol{\beta}_t$ throughout the computation graph allows every downstream module to form its own regime-weighted expectation, gracefully degrading toward more conservative behaviour as regime uncertainty grows.

\section{Regime State Representation}

\begin{definition}[MarketRegimeState]
The \texttt{MarketRegimeState} at time $t$ is a tuple:
\begin{equation}
    \mathcal{R}_t = \left(
        \boldsymbol{\beta}_t,\; H_t,\; \hat{U}_t,\; \mathbf{f}_t,\; v_t
    \right),
\end{equation}
where:
\begin{enumerate}
    \item $\boldsymbol{\beta}_t \in \Delta^{K-1}$ is the posterior probability
          vector over $K$ regime labels, lying in the $(K-1)$-dimensional
          probability simplex so that $\beta_t(k) \geq 0$ for all $k$ and
          $\sum_{k=1}^{K} \beta_t(k) = 1$.
    \item $H_t = -\sum_{k=1}^K \beta_t(k) \ln \beta_t(k)$ is the Shannon
          entropy of the regime posterior, attaining its minimum value of zero
          when the classifier is certain and its maximum value of $\ln K$ when
          all regimes are equiprobable.
    \item $\hat{U}_t = H_t / \ln K \in [0, 1]$ is the normalised uncertainty
          index, scaled so that $\hat{U}_t = 0$ denotes complete certainty and
          $\hat{U}_t = 1$ denotes maximal ignorance regardless of the number of
          regime labels $K$.
    \item $\mathbf{f}_t \in \mathbb{R}^d$ is the contributing feature vector
          consumed by the regime classifier at time $t$, stored verbatim for
          post-hoc attribution, debugging, and model monitoring purposes.
    \item $v_t \in \mathbb{S}$ is the model version string of the regime
          estimator that produced $\boldsymbol{\beta}_t$, enabling downstream
          consumers to condition their behaviour on known model-specific
          biases or calibration epochs.
\end{enumerate}
\end{definition}

The five-component structure of $\mathcal{R}_t$ is deliberate and reflects hard lessons learned from prior systems that transmitted only coarse regime labels across component boundaries. Storing only the MAP estimate $\hat{R}_t = \arg\max_k \beta_t(k)$ would discard the entire posterior's shape, collapsing a rich probability distribution over $K$ regimes into a single integer label and making it impossible to distinguish between a confident classification---where one regime accounts for, say, $95\%$ of the posterior mass---and a near-uniform distribution in which no regime commands more than $1/K$ of the probability, two situations that demand qualitatively different downstream responses in terms of position sizing, stop-loss calibration, and signal aggregation. The Shannon entropy $H_t$ and its normalised form $\hat{U}_t$ make this distinction explicit and machine-readable, allowing downstream operators to implement smooth, entropy-weighted attenuation of risk budgets rather than hard threshold logic. This continuous degradation of confidence translates directly into proportional reductions in allocated notional exposure, preventing the sharp, discontinuous changes in portfolio composition that would otherwise occur at entropy thresholds and that can themselves generate adverse market impact.

Including the raw feature vector $\mathbf{f}_t$ as part of the state record serves a purpose that goes beyond simple logging. Given $\mathcal{R}_t$ and the model identified by $v_t$, the regime posterior $\boldsymbol{\beta}_t$ can be exactly reproduced offline---to numerical precision, subject only to floating-point determinism---without access to any other runtime state, including the live data feed, the feature engineering pipeline, or the regime model's internal state buffer. This self-contained reproducibility property is essential for at least four distinct operational requirements. First, it satisfies the regulatory auditability demands of frameworks such as MiFID~II and SEC Rule~15c3-5, which require firms to demonstrate that pre-trade risk checks were active and correctly parameterised at the moment each order was submitted; the regime posterior is a direct input to those pre-trade checks, and its exact reproducibility closes the evidentiary chain. Second, it enables post-incident debugging of anomalous trading behaviour: when a position unexpectedly grows to an outsized size or a stop-loss fires at an unanticipated level, the stored $\mathcal{R}_t$ record allows engineers to replay the exact decision logic that was operative at that moment, isolating whether the anomaly originated in the feature engineering layer, the regime estimation layer, or the downstream signal or risk layer. Third, it supports continuous model validation processes that compare live regime outputs against retrospective ground-truth regime labels derived from realised return statistics, volatility clustering, and drawdown analysis; because the feature vector is stored, these comparisons can be made at the feature level as well as the label level, enabling attribution of performance degradation to specific feature groups. Fourth, it facilitates A/B testing of new model versions in shadow mode, because the shadow model can be applied to the exact same feature vector that drove the production model's output without requiring the shadow system to independently reconstruct the feature state.

The model version string $v_t \in \mathbb{S}$ encodes not just a semantic version number but a complete, content-addressed identifier---implemented in production as a SHA-256 hash of the serialised model artefact, including all parameter tensors and preprocessing pipeline specifications---so that the exact parameter set that produced $\boldsymbol{\beta}_t$ can be retrieved from the model registry without ambiguity. This is particularly important when multiple model versions are simultaneously active across different deployment environments, as can occur during rolling upgrades or when different asset classes are served by regime models that were retrained on different data vintages. The version string ensures that consumers of $\mathcal{R}_t$ can condition their behaviour on known model-specific biases or calibration epochs: a consumer that has characterised the systematic over-confidence of model version $v_a$ relative to $v_b$ in high-volatility environments can apply a version-conditioned posterior recalibration before forming its own decisions, without needing to modify the regime engine itself.

The MAP regime estimate is formally:
\begin{equation}
    \hat{R}_t = \arg\max_{k \in \{1,\ldots,K\}} \beta_t(k),
\end{equation}
and it is retained in the state record as a convenience accessor for monitoring dashboards and human-readable logging, but it plays no direct role in the quantitative decision pipeline. Critically, $\hat{R}_t$ alone is never passed downstream as the sole regime signal. The full posterior vector $\boldsymbol{\beta}_t \in \Delta^{K-1}$ is always propagated so that downstream operators can form regime-weighted expectations of any quantity they compute. The principle underlying this design is that the regime posterior is a sufficient statistic for all downstream decisions that depend on regime identity: any function of the regime label that a downstream operator might wish to compute can be expressed as an expectation under $\boldsymbol{\beta}_t$, and discarding $\boldsymbol{\beta}_t$ in favour of $\hat{R}_t$ constitutes an information loss that cannot be recovered. For instance, a signal generator computing the expected return $\mu_t$ under regime uncertainty would compute the regime-weighted blend
\begin{equation}
    \mu_t = \sum_{k=1}^{K} \beta_t(k)\, \mu_t^{(k)},
\end{equation}
where $\mu_t^{(k)}$ is the expected return conditional on regime $k$ being active, estimated from the regime-conditional return history. The variance of $\mu_t$ induced purely by regime uncertainty---distinct from the within-regime return variance---is
\begin{equation}
    \mathrm{Var}_{\boldsymbol{\beta}_t}[\mu_t^{(\cdot)}] = \sum_{k=1}^{K} \beta_t(k)\, \bigl(\mu_t^{(k)} - \mu_t\bigr)^2,
\end{equation}
a regime-uncertainty risk premium that should be reflected in position sizing even when the within-regime return forecasts are individually precise. A portfolio optimiser that operates on hard regime labels cannot construct this term at all, as it has no representation of the probability mass assigned to regimes other than the MAP estimate. The soft-assignment approach is therefore not merely a stylistic preference but a structural requirement for any downstream operator that needs to account for the full decision-theoretic consequences of regime ambiguity.

Beyond signal blending, the soft-assignment framework resolves a well-known and serious deficiency of hard regime switching in backtesting environments. When a hard classifier changes its output from regime $r_i$ to regime $r_j$, the precise calendar date at which this transition occurred is retrospectively identifiable with far greater precision than is possible in real time: a backtester can observe the full realised feature trajectory and apply the classifier's threshold function to determine, to the exact day, when the switch would have occurred. In contrast, the real-time observer at time $t$ holds only the filtered posterior $\boldsymbol{\beta}_t$, which may assign substantial probability to both $r_i$ and $r_j$ simultaneously for many consecutive periods before the posterior concentrates. A backtest that uses hard regime labels therefore implicitly grants itself foreknowledge of transition timing that was not available to the live system, introducing a form of look-ahead bias that inflates the measured performance of regime-conditional strategies. The soft-assignment framework eliminates this asymmetry by propagating the same posterior object in both live operation and backtesting, so that the backtest's regime information content exactly matches what was available to the live system at each point in time, subject only to the constraint that the model's parameters must have been estimated on data strictly prior to the backtest period.




\section{Hidden Markov Model Foundation}

The regime sequence $\{R_t\}_{t \geq 0}$ is modelled as a discrete-state, first-order Markov chain with finite state space $\mathcal{S} = \{r_1, r_2, \ldots, r_K\}$ and row-stochastic transition matrix $\Pi \in \mathbb{R}^{K \times K}$:
\begin{equation}
    P(R_{t+1} = r_j \mid R_t = r_i) = \Pi_{ij}, \qquad \sum_{j=1}^K \Pi_{ij} = 1 \quad \forall i.
\end{equation}
The Markov property encodes the economically motivated assumption that the regime at time $t+1$ depends on history only through the current regime $R_t$, not through the entire trajectory $\{R_s\}_{s < t}$. This is a modelling simplification; longer-range dependencies can be partially accommodated by increasing $K$ or by augmenting the state space with lagged regime indicators, but empirical studies on monthly equity data suggest that the first-order approximation captures the majority of regime persistence structure at the horizons relevant to daily trading.

The transition matrix $\Pi$ encodes the persistence and typical duration of each regime. The expected sojourn time in regime $r_i$ before a transition occurs is $\mathbb{E}[\tau_i] = (1 - \Pi_{ii})^{-1}$, a quantity that is directly observable from historical data and serves as a sanity check during model validation: a bull-market regime in equity indices should have $\mathbb{E}[\tau_i]$ measured in months to years, while a crisis regime should have a shorter expected duration but a heavier tail, reflecting the possibility of prolonged distress periods such as those observed during the 2008--2009 global financial crisis or the 2020 pandemic shock. Off-diagonal elements $\Pi_{ij}$ for $i \neq j$ carry the probability of transitioning from regime $r_i$ to regime $r_j$ in a single time step; in practice, many of these entries are estimated to be very small, as direct transitions between structurally distant regimes (e.g., from a low-volatility trending regime directly to a crisis regime, bypassing an intermediate elevated-volatility regime) are relatively rare.

The observation model assumes that the feature vector $\mathbf{x}_t \in \mathbb{R}^d$ is drawn from a regime-conditioned multivariate Gaussian:
\begin{equation}
    \mathbf{x}_t \mid R_t = r_k \;\sim\; \mathcal{N}(\boldsymbol{\mu}_k, \boldsymbol{\Sigma}_k),
\end{equation}
where $\boldsymbol{\mu}_k \in \mathbb{R}^d$ and $\boldsymbol{\Sigma}_k \in \mathbb{S}_{++}^d$ are the regime-specific mean vector and positive definite covariance matrix. The Gaussian assumption is a pragmatic choice that admits closed-form evaluation of the likelihood $p(\mathbf{x}_t \mid R_t = r_k)$ and lends itself to efficient Baum-Welch parameter estimation via the expectation-maximisation algorithm. In practice, the tails of financial return distributions are heavier than Gaussian, and regime-conditioned Student-$t$ observation models have been studied as alternatives; however, for the feature vectors used here---which include volatility estimates, momentum signals, and macro indicators rather than raw returns---the Gaussian approximation is substantially more defensible, and robustness to tail misspecification is further mitigated by the ensemble architecture described in Section~\ref{sec:regime_ensemble}.

The real-time filtered belief update at each new observation $\mathbf{x}_t$ follows the standard HMM forward algorithm, which applies Bayes' rule conditioned on the Markov transition:
\begin{equation}
    \beta_t(k) = \frac{p(\mathbf{x}_t \mid R_t = r_k)\, \sum_{i=1}^{K} \Pi_{ik}\, \beta_{t-1}(i)}
    {\displaystyle\sum_{j=1}^{K} p(\mathbf{x}_t \mid R_t = r_j)\, \sum_{i=1}^{K} \Pi_{ij}\, \beta_{t-1}(i)}.
    \label{eq:hmm_filter}
\end{equation}
The numerator decomposes into two interpretable factors: $\sum_{i=1}^{K} \Pi_{ik}\, \beta_{t-1}(i)$ is the one-step-ahead predictive probability of being in regime $r_k$ at time $t$, obtained by marginalising the previous posterior through the transition matrix, and $p(\mathbf{x}_t \mid R_t = r_k)$ is the likelihood of the observed feature vector under the emission model for regime $r_k$. The denominator is the total predictive likelihood, serving as the normalisation constant that ensures $\boldsymbol{\beta}_t \in \Delta^{K-1}$. This recursive update is $\mathcal{O}(K^2 d)$ per time step, making it computationally tractable even for large feature dimensions $d$ and moderately large regime counts $K \leq 12$.

Model parameters $\theta = \{\Pi, \{\boldsymbol{\mu}_k, \boldsymbol{\Sigma}_k\}_{k=1}^K\}$ are estimated offline via the Baum-Welch algorithm applied to a training corpus spanning at least one full market cycle, ideally encompassing multiple distinct volatility and trend regimes. The resulting parameter estimates are treated as fixed during live operation and refreshed on a scheduled cadence, with the model version string $v_t$ updated upon each re-estimation. A shadow deployment infrastructure runs the newly estimated model in parallel for a configurable validation period before it replaces the production model, preventing silent degradation from parameter drift.

\section{Multi-Model Regime Ensemble}
\label{sec:regime_ensemble}

No single modelling approach to regime identification is universally dominant across all market conditions, feature sets, and estimation horizons. Gaussian HMMs are well-suited to capturing slow-moving, persistent structural shifts but can be slow to respond to sudden regime transitions, particularly in crisis periods where the feature distribution shifts abruptly and the filtered belief update requires several observations before the posterior concentrates on the new regime. Threshold-based classifiers respond immediately to threshold crossings in observable stress indicators but are brittle to the specific threshold values chosen and can generate spurious regime switches during brief, transient spikes in those indicators. Machine-learning classifiers can exploit non-linear interactions among high-dimensional feature sets that are invisible to parametric models, but they require careful regularisation to avoid overfitting to historical crisis episodes that may not repeat in precisely the same form. Cross-asset correlation regime detectors capture systemic risk propagation dynamics that are invisible to single-asset feature sets, but they require sufficient cross-sectional breadth to be statistically reliable and are inherently lagging due to the rolling estimation windows required. For these reasons, the platform runs four complementary regime estimators in parallel, each contributing an independent posterior distribution over the same $K$-label regime space:

\begin{enumerate}
    \item \textbf{Gaussian Mixture HMM.} A Hidden Markov Model with Gaussian mixture emission densities, fitted on a $d_1$-dimensional price-based feature vector encompassing realised volatility at multiple horizons (5-day, 21-day, 63-day), signed momentum at those same horizons, volume-weighted average price deviation from moving averages, and the ratio of upward to downward daily closes over rolling windows. The mixture emission model, with $M$ components per regime, relaxes the unimodal Gaussian constraint and accommodates the multimodal feature distributions that arise when a single economic regime label subsumes structurally distinct sub-periods.

    \item \textbf{Threshold-Based Macro Stress Classifier.} A rule-based regime classifier that maps observable macro-financial stress indicators to regime labels via learned threshold functions. The primary indicators include the CBOE Volatility Index (VIX) and its term structure (VIX3M/VIX ratio), investment-grade and high-yield credit spreads, the slope and curvature of the US Treasury yield curve, and the TED spread as a measure of interbank funding stress. Thresholds are calibrated on historical data using isotonic regression to respect the natural ordering of stress levels across regime labels. This classifier has zero latency relative to the HMM's multi-step belief convergence and serves as an early-warning layer in crisis onset scenarios.

    \item \textbf{Gradient-Boosted Macro Classifier.} A supervised gradient boosting classifier (XGBoost with depth-limited trees and $\ell_2$ regularisation) trained on a $d_3$-dimensional macroeconomic feature vector that includes industrial production growth, PMI survey readings, leading economic index components, consumer confidence indices, housing starts, and labour market indicators. The training labels are derived from retrospective regime annotations produced by a panel of domain experts cross-validated against quantitative regime clustering on realised return data. Probabilistic outputs are obtained via Platt scaling applied to the raw classifier scores, yielding calibrated posterior probabilities that are directly comparable to the HMM and threshold classifier outputs.

    \item \textbf{Rolling-Window Correlation Regime Detector.} A structural break detector that operates on the $(n \times n)$ cross-asset realised correlation matrix $\mathbf{C}_t$, estimated over a rolling window of $\tau$ trading days using the DCC-GARCH framework to control for heteroskedasticity-induced correlation distortion. Regime transitions are detected by monitoring the Frobenius-norm distance $\|\mathbf{C}_t - \mathbf{C}_{t-\tau}\|_F$ and the leading eigenvalue $\lambda_1(\mathbf{C}_t)$, which measures the degree of synchronisation across assets and is known to spike sharply during systemic risk episodes. This detector is particularly sensitive to the onset of correlation breakdowns---periods during which historically low-correlated assets suddenly co-move, eliminating the diversification benefits assumed by portfolio construction models calibrated in normal regimes.
\end{enumerate}

The four estimators produce independent posterior distributions $\boldsymbol{\beta}_t^{(m)}$ for $m \in \{1, 2, 3, 4\}$ over the shared $K$-label regime space. Their outputs are fused into an ensemble posterior via a learned convex combination:
\begin{equation}
    \boldsymbol{\beta}_t^{\mathrm{ens}} = \sum_{m=1}^{4} w_m\, \boldsymbol{\beta}_t^{(m)}, \qquad w_m \geq 0,\quad \sum_{m=1}^{4} w_m = 1,
\end{equation}
where the mixture weights $\{w_m\}$ are estimated by maximising the log predictive score of the ensemble on a held-out validation set. The disagreement between estimators is quantified as the average pairwise Jensen-Shannon divergence:
\begin{equation}
    D_t^{\mathrm{JSD}} = \frac{2}{M(M-1)} \sum_{m < m'} \mathrm{JSD}\!\left(\boldsymbol{\beta}_t^{(m)} \,\|\, \boldsymbol{\beta}_t^{(m')}\right),
\end{equation}
where $\mathrm{JSD}(P \| Q) = \frac{1}{2} D_{\mathrm{KL}}(P \| M) + \frac{1}{2} D_{\mathrm{KL}}(Q \| M)$ with $M = \frac{1}{2}(P + Q)$, and $D_{\mathrm{KL}}$ denotes the Kullback-Leibler divergence. The Jensen-Shannon divergence is chosen over the raw KL divergence because it is symmetric, always finite (even when one distribution assigns zero probability to a label that another assigns positive probability), and bounded in $[0, \ln 2]$, making it amenable to threshold-based flagging. When $D_t^{\mathrm{JSD}}$ exceeds a configurable threshold $\delta_{\mathrm{JSD}}$, the normalised uncertainty flag $\hat{U}_t$ is elevated irrespective of the entropy of the ensemble posterior itself, and all downstream operators receiving $\mathcal{R}_t$ apply more conservative position sizing, widen their stop-loss tolerances, and increase their required signal-to-noise thresholds before committing capital.

\section{Regime Features}

\begin{table}[H]
\centering
\caption{Feature categories used by the market regime engine.}
\label{tab:regime_features}
\small
\begin{tabularx}{\textwidth}{p{4.2cm} p{10.3cm}}
\toprule
\textbf{Feature Category} & \textbf{Examples} \\
\midrule
Trend indicators & 20/50/200-day moving average crossovers, ADX,
    directional movement index. \\
\midrule
Volatility indicators & Realised volatility (1-day to 1-month),
    implied volatility (VIX, VVIX), GARCH conditional volatility. \\
\midrule
Liquidity indicators & Bid-ask spreads, Amihud illiquidity ratio,
    turnover relative to 90-day average. \\
\midrule
Cross-asset correlation & Rolling pairwise correlations between equities,
    rates, credit, commodities, and currencies. \\
\midrule
Macroeconomic indicators & Yield curve slope (10y--2y spread), credit
    spread (IG OAS), breakeven inflation, PMI, unemployment. \\
\midrule
Market microstructure & Order flow imbalance, quote stuffing frequency,
    dark pool volume ratio, short interest. \\
\midrule
Dispersion indicators & Cross-sectional return dispersion, sector
    rotation speed, factor loading stability. \\
\bottomrule
\end{tabularx}
\end{table}

% ============================================================
% CHAPTER 6: RISK MANAGEMENT FRAMEWORK
% ============================================================
\chapter{Risk Management Framework}

\section{The Risk-First Principle}

The most dangerous failure mode in algorithmic trading is not a bad
prediction; it is a good prediction that is acted upon without proper
risk controls. The platform enforces the risk-first principle as an
architectural invariant: every proposed action must pass through a
\texttt{RiskOperator} before it reaches any execution layer. The
\texttt{RiskOperator} is not advisory; it has veto power.

\section{Position Sizing: Kelly Criterion}

Given the win probability $p$ of a strategy, the loss probability
$q = 1 - p$, and the win/loss payoff ratio $b$, the Kelly-optimal
fraction of capital to allocate is:
\begin{equation}
    f^* = p - \frac{q}{b} = \frac{pb - q}{b}.
    \label{eq:kelly}
\end{equation}
The second-order condition confirms this is a maximum:
\begin{equation}
    \frac{d^2G}{df^2} = -\frac{pb^2}{(1+bf)^2} - \frac{q}{(1-f)^2} < 0,
\end{equation}
where $G(f) = p \ln(1 + bf) + q\ln(1-f)$ is the expected logarithmic
growth rate. The platform applies a fractional Kelly with parameter
$\kappa \in (0,1]$:
\begin{equation}
    f = \kappa f^*,
\end{equation}
and always imposes a hard cap $f \leq f_{\max}$ regardless of the Kelly
output, as a guard against estimation error in $p$ and $b$.

\section{Risk Metrics}

\subparagraph{Sharpe Ratio.}
\begin{equation}
    \Sharpe = \frac{\mathbb{E}[R_p - R_f]}{\sigma_p} \cdot \sqrt{T},
\end{equation}
where $R_p$ is portfolio return, $R_f$ is the risk-free rate, $\sigma_p$
is annualised standard deviation, and $T$ is the number of periods per year.

\subparagraph{Sortino Ratio.}
\begin{equation}
    \text{Sortino} = \frac{\mathbb{E}[R_p - R_f]}{\sigma_d} \cdot \sqrt{T},
    \qquad
    \sigma_d = \sqrt{\mathbb{E}\!\left[\min(R_p - R_{\mathrm{tgt}}, 0)^2\right]}.
\end{equation}

\subparagraph{Conditional Value-at-Risk.}
For portfolio loss $L = -R_p$ at confidence level $\alpha \in (0,1)$:
\begin{equation}
    \VaR_\alpha = \inf\{l : P(L > l) \leq 1 - \alpha\},
    \qquad
    \CVaR_\alpha = \mathbb{E}[L \mid L \geq \VaR_\alpha].
\end{equation}
The linear programming formulation over $T$ empirical scenarios is:
\begin{equation}
    \CVaR_\alpha = \min_{z \in \mathbb{R}}
    \left\{
        z + \frac{1}{(1-\alpha)T} \sum_{t=1}^T \max(\ell_t - z, 0)
    \right\}.
    \label{eq:cvar_lp}
\end{equation}

\subparagraph{Maximum Drawdown.}
\begin{equation}
    \text{MaxDD} = \max_{t} \frac{\max_{s \leq t} W_s - W_t}{\max_{s \leq t} W_s},
\end{equation}
where $W_t$ is portfolio wealth at time $t$.

\section{Risk Dimensions and Limits}

\begin{table}[H]
\centering
\caption{Risk dimensions monitored by the \texttt{RiskOperator}.}
\label{tab:risk_dimensions}
\small
\begin{tabularx}{\textwidth}{p{4.0cm} p{10.5cm}}
\toprule
\textbf{Risk Dimension} & \textbf{Measurement and Action} \\
\midrule
Market risk &
    Portfolio delta, beta, and factor exposures. Excess exposure triggers
    hedging request to \texttt{HedgeOptimizationOperator}. \\
\midrule
Volatility risk &
    Ratio of current realised volatility to 90-day average. Position sizes
    scaled inversely with volatility. \\
\midrule
Liquidity risk &
    Order size relative to average daily volume. Orders exceeding 1\% of ADV
    are split or rejected. \\
\midrule
Concentration risk &
    Single-name, sector, and factor concentration limits. Breaches trigger
    forced diversification. \\
\midrule
Correlation risk &
    Rolling pairwise correlation between held positions. High correlation
    reduces effective diversification benefit. \\
\midrule
Leverage &
    Gross and net leverage relative to NAV. Hard ceiling enforced. \\
\midrule
Drawdown &
    Current drawdown relative to high-water mark. Breaching soft limit
    reduces position sizes; breaching hard limit halts new positions. \\
\midrule
Tail risk (CVaR) &
    CVaR at 95\% confidence. Breaching limit halts new position-opening
    until resolved. \\
\midrule
Model risk &
    Regime uncertainty $\hat{U}_t$. When $\hat{U}_t > \theta_U$,
    conservative sizing applies. \\
\midrule
Stale data risk &
    Any input datum flagged \texttt{STALE} or \texttt{SUSPECT} triggers
    a conservative fallback or order pause. \\
\bottomrule
\end{tabularx}
\end{table}

\section{Minimum-Variance Hedge Ratio}

The minimum-variance hedge ratio between a portfolio with return $R_V$
and a hedge instrument with return $R_F$ is:
\begin{equation}
    h^* = \frac{\mathrm{Cov}(R_V, R_F)}{\mathrm{Var}(R_F)}
        = \rho_{VF} \cdot \frac{\sigma_V}{\sigma_F},
    \label{eq:hedge_ratio}
\end{equation}
with hedge effectiveness:
\begin{equation}
    \eta = 1 - \frac{\mathrm{Var}(R_H)}{\mathrm{Var}(R_V)} = \rho_{VF}^2.
\end{equation}
The \texttt{HedgeOptimizationOperator} recomputes $h^*$ dynamically
using regime-conditioned correlation estimates from the
\texttt{MarketRegimeOperator}.

\section{Risk-Adjusted Objective Function}

The platform optimises the risk-adjusted cumulative return over
horizon $T$:
\begin{equation}
    \mathcal{J}(\pi) = \mathbb{E}\!\left[\sum_{t=0}^T \gamma^t r_t\right]
    - \lambda_1 \cdot \CVaR_\alpha(\pi)
    - \lambda_2 \cdot \text{MaxDD}(\pi),
    \label{eq:objective}
\end{equation}
where $\gamma \in (0, 1]$ is the discount factor and $\lambda_1, \lambda_2 \geq 0$
are penalty weights calibrated to the operator's risk tolerance.

% ============================================================
% CHAPTER 7: EXECUTION AND SIMULATION LAYERS
% ============================================================
\chapter{Execution and Simulation Layers}

\section{Simulation-First Architecture}

The platform is designed from the outset to operate in simulation mode.
All execution components default to paper trading. The interface hierarchy
is:

\begin{table}[H]
\centering
\caption{Execution interface hierarchy.}
\label{tab:exec_hierarchy}
\small
\begin{tabularx}{\textwidth}{p{4.5cm} p{10.0cm}}
\toprule
\textbf{Interface} & \textbf{Description} \\
\midrule
\texttt{ExecutionSimulator} &
    Historical replay mode. Orders are filled at OHLCV prices with
    configurable slippage and transaction cost models. \\
\midrule
\texttt{PaperBroker} &
    Real-time paper trading mode. Orders are submitted to a
    simulated order book with realistic latency and partial-fill modelling. \\
\midrule
\texttt{LiveBrokerAdapter} &
    Live execution adapter, architecturally separated and deliberately
    absent from the initial implementation. Requires explicit configuration
    and environment flag to activate. \\
\bottomrule
\end{tabularx}
\end{table}

\section{Transaction Cost Model}

Naive backtesting that ignores transaction costs systematically overstates strategy performance, often by margins large enough to render an apparently profitable strategy unprofitable in live deployment. The failure mode is well-documented in the academic literature and stems from the implicit assumption that every trade executes at the quoted mid-price with zero friction---an assumption that is never satisfied in practice. To prevent this class of overfitting, the platform implements a layered, composable transaction cost model in which every simulated order is subjected to four distinct cost components, each modelled separately and applied in sequence before the resulting net P\&L is recorded.

\textbf{Commission} represents the explicit, contractual cost charged by the executing broker for each completed trade. This cost is configurable at the asset-class level and further parameterised by broker tier, reflecting the reality that institutional clients operating under prime brokerage agreements face materially lower per-share or per-contract fees than retail participants. Commissions are expressed in basis points of notional value for equity and fixed-income instruments, and in per-contract dollar terms for listed derivatives. The configuration is loaded at initialisation time from a typed \texttt{CommissionSchedule} object, ensuring that any downstream change to broker terms is immediately and consistently propagated to all simulated execution paths.

\textbf{Bid-ask spread} captures the implicit cost of crossing the market. Rather than treating the spread as a fixed constant---a common but inaccurate simplification---the platform estimates the half-spread dynamically from a rolling microstructure model calibrated to recent order-book data. Specifically, the half-spread $s_t$ at time $t$ is estimated as a weighted average of the Roll implied spread estimator \citep{Roll1984} and the Corwin--Schultz high-low estimator, with the weighting determined by the availability and quality of tick-level data. For each simulated order, the fill price is penalised by $s_t$ relative to the mid-price, so that a buy order executes at the ask side and a sell order executes at the bid side of the estimated spread. This approach captures the time-varying nature of liquidity costs across different market regimes and intraday periods, since spreads widen significantly during the open and close auctions, around macroeconomic announcements, and during stress episodes.

\textbf{Market impact} accounts for the adverse price movement caused by the act of executing a large order relative to the prevailing liquidity. The platform models temporary market impact using the Almgren--Chriss square-root law \citep{AlmgrenChriss2001}, which has strong empirical support across asset classes and is derived from a continuous-time optimal execution framework. The impact cost assigned to an order of size $V$ executed against an instrument with average daily volume $\text{ADV}$ and daily volatility $\sigma$ is:
\begin{equation}
    \text{Impact}(V) = \gamma \cdot \sigma \cdot \sqrt{\frac{V}{\text{ADV}}},
\end{equation}
where $\gamma > 0$ is a dimensionless liquidity coefficient that is calibrated instrument-by-instrument from historical execution data or, in the absence of proprietary execution records, sourced from published empirical estimates in the market microstructure literature. The square-root functional form reflects the concavity of the price-impact schedule: doubling the order size increases costs by a factor of $\sqrt{2}$, not $2$, because larger orders tend to attract more passive liquidity as they work through the book. Permanent impact---the component of price displacement that persists after the order is complete---is modelled separately as a linear fraction $\eta \cdot V / \text{ADV}$ and is applied to the mark-to-market price of the position for all subsequent periods, correctly capturing the information signal that large trades convey to the broader market.

\textbf{Slippage} models the residual timing cost that arises from the latency between a signal being generated and the order being filled. Even after accounting for the bid-ask spread and market impact, a realised fill price will typically differ from the signal price due to intraday price drift during the execution window. The platform models this component by drawing timing slippage realisations from the empirical distribution of intraday price moves observed during the execution window, conditioned on the prevailing volatility regime and time-of-day. Concretely, for each simulated order, a slippage draw $\varepsilon \sim \hat{F}_{\text{slip}}(\cdot \mid \sigma_t, \text{tod})$ is sampled from the regime-conditioned empirical distribution estimated over the trailing lookback window, and the fill price is further adjusted by $\varepsilon$. This approach captures the asymmetric and fat-tailed nature of timing slippage that is missed by Gaussian approximations, and ensures that the backtesting engine does not implicitly assume instantaneous, costless execution.

The total round-trip transaction cost for a trade of size $V$ is therefore the sum of all four components, computed independently for the entry and exit legs. The \texttt{TransactionCostEngine} exposes a typed interface that accepts an \texttt{OrderIntent} and returns a \texttt{CostBreakdown} struct, making the individual cost components auditable and allowing research analysts to isolate the contribution of each layer to overall strategy drag. Strategies that remain profitable after subtracting the full layered cost model are substantially more likely to replicate that performance in live execution.

\section{Walk-Forward Backtesting}

Walk-forward validation is the methodological standard for evaluating predictive models in financial time series, and the platform enforces it as the default validation protocol for all strategy evaluation pipelines. The core motivation is straightforward: financial return data exhibit non-stationarity, structural breaks, and slow-moving regime changes that render cross-validation schemes designed for i.i.d. data---such as $k$-fold or shuffle-split---fundamentally invalid. A model that is evaluated on data that temporally precedes its training set has implicitly been granted access to future information, a condition known as look-ahead bias, which produces wildly optimistic performance estimates that cannot be reproduced in deployment. Walk-forward validation eliminates this failure mode by enforcing strict temporal ordering: the model is always tested on data that lies entirely in the future of its training set.

Formally, let $\mathcal{D} = \{(x_t, y_t)\}_{t=1}^{T}$ denote the full dataset ordered chronologically. The walk-forward scheme is parameterised by an in-sample window size $W$, a step size $S$, a purging window $P$, and an embargo $E$. Beginning at anchor $t_0$, the procedure iterates as follows. In the first step, the model is trained on the in-sample window $[t_0,\, t_0 + W)$, which contains $W$ consecutive observations. The training objective, regularisation hyperparameters, and feature selection are all determined exclusively from this window, with no reference to any data outside it. In the second step, a purging window of length $P$ immediately following the training set is excluded from both training and testing. The purpose of purging is to eliminate observations whose labels $y_t$ overlap in time with the training labels, a subtle but critical source of leakage identified by \citet{Lopez2018} that arises whenever the target variable is computed over a multi-period horizon, such as a five-day forward return. In the third step, an embargo of length $E$ is additionally excluded immediately after the purging window. The embargo guards against the residual autocorrelation structure of financial returns causing the model to effectively ``see'' features from the training period through their lagged influence on the test period. Only after both the purging and embargo windows have been skipped does the test window begin, spanning $[t_0 + W + P + E,\, t_0 + W + P + E + S)$. Performance metrics are recorded for this test window, and the anchor is then advanced by $t_0 \leftarrow t_0 + S$. The entire procedure repeats until the test window would exceed the end of the dataset.

The resulting sequence of out-of-sample performance observations $\{\widehat{\text{SR}}_k\}_{k=1}^{K}$ constitutes the walk-forward performance profile, where $K = \lfloor (T - W - P - E) / S \rfloor$ is the number of folds. The platform reports the time-averaged Sharpe ratio, its standard error across folds, the fraction of folds with positive returns, and the worst-fold drawdown, providing a distributional view of strategy robustness rather than a single point estimate. Strategies with high mean Sharpe but large inter-fold variance are flagged as regime-sensitive and subjected to additional stress testing before live deployment approval.

The purging and embargo windows are not optional configuration parameters; they are enforced programmatically by the \texttt{ValidationOperator} at each fold boundary. Any attempt to instantiate a walk-forward experiment without specifying non-negative values for $P$ and $E$ raises a \texttt{ValidationConfigurationError} at initialisation time. The default values are $P = 5$ trading days and $E = 2$ trading days, reflecting a five-day forward return horizon and a two-day autocorrelation decay buffer, but these defaults are overridden automatically when the \texttt{StrategyManifest} declares a different label horizon. This tight coupling between the strategy definition and the validation protocol ensures that the appropriate anti-leakage windows are always applied without requiring the researcher to manually compute and set them.

\section{Monte Carlo Simulation}

The Monte Carlo simulation engine provides a complementary stress-testing capability to the walk-forward backtester, addressing the fundamental limitation that historical data contains only one realisation of the market process. No matter how carefully a walk-forward experiment is conducted, it evaluates the strategy against a single historical path, which may be atypically benign or atypically hostile in ways that are difficult to diagnose. Monte Carlo simulation addresses this by generating a large ensemble of synthetic price paths that are statistically consistent with the estimated data-generating process but span a much wider range of market conditions than any single historical record, enabling the analyst to form a distributional view of strategy performance under genuine uncertainty about future market dynamics.

The engine operates in four stages. In the first stage, the regime-conditioned return distribution is estimated from the filtered posterior $\boldsymbol{\beta}_t$ produced by the Hidden Markov Model described in Chapter~5. Conditional on the current regime $r \in \{1, \ldots, K\}$, the factor return vector $\mathbf{f}_t \in \mathbb{R}^d$ is modelled as a multivariate location-scale distribution with regime-specific mean $\boldsymbol{\mu}^{(r)}$ and covariance $\boldsymbol{\Sigma}^{(r)}$. The regime probabilities at the simulation start date are drawn from the filtered posterior, so that the ensemble of paths correctly reflects the current uncertainty about which regime the market is operating in, rather than conditioning on a single point estimate of the regime state.

In the second stage, correlated factor shocks are drawn jointly from the regime-conditioned covariance matrix $\boldsymbol{\Sigma}^{(r)}$ using a Cholesky decomposition. Let $\boldsymbol{\Sigma}^{(r)} = \mathbf{L}^{(r)} (\mathbf{L}^{(r)})^\top$ be the Cholesky factorisation. Each factor shock vector is constructed as $\mathbf{f}_t^{(r)} = \boldsymbol{\mu}^{(r)} + \mathbf{L}^{(r)} \mathbf{z}_t$, where $\mathbf{z}_t$ is a vector of i.i.d.\ standardised draws. This construction ensures that the cross-sectional dependence structure encoded in $\boldsymbol{\Sigma}^{(r)}$ is preserved across all simulated paths, which is essential for correctly capturing diversification effects and tail dependence during stress scenarios.

In the third stage, fat-tail corrections are applied by drawing the standardised innovations $\mathbf{z}_t$ from a multivariate Student-$t$ distribution rather than from a multivariate Gaussian. The degrees-of-freedom parameter $\nu^{(r)}$ is calibrated separately for each regime by maximum likelihood estimation on the historical residuals after removing the Gaussian mean and covariance structure. Empirically, $\nu^{(r)}$ ranges from approximately $4$ in high-volatility regimes---where tail events are frequent---to $8$--$12$ in low-volatility regimes where the return distribution is closer to Gaussian. The Student-$t$ residuals introduce the excess kurtosis necessary to generate the occasional large-magnitude shocks that characterise financial markets but are systematically underrepresented in Gaussian simulations, ensuring that risk measures derived from the simulation are not overly optimistic about tail behaviour.

In the fourth stage, $N = 10{,}000$ independent paths are generated by applying the above procedure sequentially over the simulation horizon $H$, with regime transitions drawn from the estimated HMM transition matrix $\mathbf{A}$ at each step. The resulting ensemble of $N$ P\&L trajectories is used to compute the empirical CVaR at confidence levels $\alpha \in \{0.95, 0.99\}$ directly from the sorted order statistics of the terminal P\&L distribution, without parametric distributional assumptions. In addition to CVaR, the engine computes the full distribution of maximum drawdown across paths, the probability of ruin (defined as the probability that the portfolio NAV falls below a configurable floor), and the expected number of margin calls under standard initial and maintenance margin rules. These outputs are surfaced through the \texttt{StressTestReport} data class and are required inputs to the pre-deployment risk approval workflow.

% ============================================================
% CHAPTER 8: POLYGLOT IMPLEMENTATION ARCHITECTURE
% ============================================================
\chapter{Polyglot Implementation Architecture}

\section{Language Role Assignments}

The platform is implemented across four languages, each assigned a
defensible and distinct architectural role.

\begin{table}[H]
\centering
\caption{Language role assignments and their justifications.}
\label{tab:lang_roles}
\small
\begin{tabularx}{\textwidth}{p{1.8cm} p{3.8cm} p{9.0cm}}
\toprule
\textbf{Language} & \textbf{Primary Role} & \textbf{Justification} \\
\midrule
Python &
    Research, orchestration, data science, ML inference &
    Richest quantitative ecosystem (NumPy, pandas, scikit-learn, PyTorch).
    Rapid strategy development and notebook-based research. Operator
    orchestration via typed Python protocols. \\
\midrule
C++ &
    Numerical kernels, pricing engines, high-throughput simulation &
    Deterministic performance for latency-sensitive computations.
    Mature linear algebra libraries (Eigen, BLAS). Options pricing,
    covariance estimation, portfolio optimisation solvers. \\
\midrule
Rust &
    Risk-critical infrastructure, concurrent event processing,
    execution simulation &
    Memory safety without garbage collection. Appropriate for the
    execution simulation engine, order management, and any component
    where correctness under concurrency is paramount. \\
\midrule
Go &
    Network services, APIs, telemetry, operational tooling &
    Excellent concurrency primitives, fast compilation, strong standard
    library for network services. Appropriate for data ingestion
    microservices, REST/gRPC APIs, and telemetry aggregation. \\
\bottomrule
\end{tabularx}
\end{table}

\section{Shared Schema Layer}

Because the platform spans four languages, shared data types must be defined
once and generated for each language. The schema layer uses Protocol Buffers
as the canonical definition format. Generated code is committed to the
repository alongside the schema definitions. Manual duplication of shared
type definitions across languages is prohibited.

Core shared types include:

\begin{table}[H]
\centering
\caption{Core shared schema types.}
\label{tab:schemas}
\small
\begin{tabularx}{\textwidth}{p{4.5cm} p{10.0cm}}
\toprule
\textbf{Type} & \textbf{Key Fields} \\
\midrule
\texttt{Signal} &
    \texttt{operator\_id}, \texttt{timestamp}, \texttt{value $\in [-1,1]$},
    \texttt{confidence}, \texttt{regime\_id}, \texttt{correlation\_id}. \\
\midrule
\texttt{OrderIntent} &
    \texttt{instrument\_id}, \texttt{side}, \texttt{quantity},
    \texttt{signal\_id}, \texttt{risk\_budget}, \texttt{urgency}. \\
\midrule
\texttt{RiskDecision} &
    \texttt{intent\_id}, \texttt{status} (APPROVED/REJECTED/RESIZED),
    \texttt{final\_quantity}, \texttt{risk\_metrics}, \texttt{reason}. \\
\midrule
\texttt{ExecutionRecord} &
    \texttt{decision\_id}, \texttt{fill\_price}, \texttt{fill\_quantity},
    \texttt{commission}, \texttt{market\_impact}, \texttt{fill\_timestamp}. \\
\midrule
\texttt{MarketRegimeState} &
    \texttt{timestamp}, \texttt{posteriors}, \texttt{entropy},
    \texttt{normalised\_uncertainty}, \texttt{model\_version}. \\
\midrule
\texttt{AuditRecord} &
    \texttt{stage}, \texttt{operator\_id}, \texttt{version},
    \texttt{input\_hash}, \texttt{output\_hash}, \texttt{duration\_ms},
    \texttt{correlation\_id}. \\
\bottomrule
\end{tabularx}
\end{table}

\section{Repository Directory Structure}

The target repository layout follows domain-driven organisation:

\begin{table}[H]
\centering
\caption{Top-level repository directory structure.}
\label{tab:repo_structure}
\small
\begin{tabularx}{\textwidth}{p{4.5cm} p{10.0cm}}
\toprule
\textbf{Directory} & \textbf{Contents} \\
\midrule
\texttt{schemas/} & Protocol Buffer definitions for all shared types. \\
\texttt{python/src/} & Python operator implementations, orchestration, research. \\
\texttt{python/tests/} & Python unit, integration, and simulation tests. \\
\texttt{cpp/include/} & C++ public headers for numerical kernels. \\
\texttt{cpp/src/} & C++ pricing engines, optimisation solvers. \\
\texttt{rust/crates/} & Rust crates for execution simulation and risk engine. \\
\texttt{go/cmd/} & Go service entry points (data ingestion, API, telemetry). \\
\texttt{go/internal/} & Go internal packages (schema adapters, middleware). \\
\texttt{strategies/} & Strategy configuration and operator composition definitions. \\
\texttt{simulations/} & Historical simulation configurations and scenario libraries. \\
\texttt{dashboards/} & Performance and risk monitoring dashboard code. \\
\texttt{docs/} & Architecture documentation, ADRs, and this whitepaper. \\
\texttt{ci/} & CI/CD pipeline definitions and quality gate scripts. \\
\texttt{config/} & Environment-specific configuration (TOML/YAML). \\
\bottomrule
\end{tabularx}
\end{table}

% ============================================================
% CHAPTER 9: TESTING AND VALIDATION TAXONOMY
% ============================================================
\chapter{Testing and Validation Taxonomy}

\section{Overview}

The platform enforces a comprehensive testing taxonomy. Each test category
targets a specific failure mode. No single test category is sufficient;
all categories are required for a merge to be approved.

\begin{table}[H]
\centering
\caption{Testing taxonomy: categories, targets, and mandatory gates.}
\label{tab:testing_taxonomy}
\small
\begin{tabularx}{\textwidth}{p{3.8cm} p{5.2cm} p{5.5cm}}
\toprule
\textbf{Test Category} & \textbf{Primary Target} & \textbf{CI Gate Requirement} \\
\midrule
Unit correctness &
    Individual operator calculation accuracy &
    100\% of operator unit tests pass. \\
\midrule
Property correctness &
    Mathematical invariants (bounds, monotonicity, idempotency) &
    All invariant assertions pass on randomised inputs. \\
\midrule
Numerical correctness &
    Floating-point accuracy relative to analytical solutions &
    Maximum relative error below $10^{-8}$ for all closed-form tests. \\
\midrule
Temporal correctness &
    Absence of forward-looking data in any pipeline stage &
    \texttt{TemporalLeakageDetector} passes with zero violations. \\
\midrule
Integration correctness &
    Correct data flow between composed operators &
    All integration tests pass in isolated test environment. \\
\midrule
Cross-language correctness &
    Numerical agreement between Python, C++, Rust, Go implementations &
    Pairwise agreement within $10^{-6}$ on shared test vectors. \\
\midrule
Strategy correctness &
    Strategy behaviour matches specification on canonical scenarios &
    All named scenario tests pass. \\
\midrule
Risk correctness &
    Risk operator correctly enforces all limit types &
    All risk boundary and veto tests pass. \\
\midrule
Simulation correctness &
    Deterministic replay of historical scenarios &
    Bit-identical output on identical seed and data. \\
\midrule
Regression correctness &
    No degradation of known-good historical cases &
    All regression snapshots pass. \\
\midrule
Performance regression &
    Operator latency within defined budgets &
    No operator exceeds its $p_{99}$ latency budget. \\
\bottomrule
\end{tabularx}
\end{table}

\section{CI/CD Pipeline Gate Sequence}

The CI/CD pipeline is structured as a strict, sequential gate sequence in which each gate must pass completely before the subsequent gate is entered. No partial advancement is permitted: a single failing assertion in any gate causes the entire pipeline to be marked failed, the corresponding pull request to be blocked from merging, and a structured failure report to be emitted to the team notification channel. This design reflects a zero-tolerance policy toward defect propagation: the cost of catching a fault at the format gate is measured in seconds, whereas the cost of the same fault reaching the performance regression gate—or, critically, a live trading session—grows by several orders of magnitude. The gate sequence is therefore both an engineering discipline and a risk-management instrument.

\begin{enumerate}
    \item \textbf{Format}: All source files must conform to language-specific,
          tool-enforced formatting standards: \texttt{black} (line length 88) for
          Python, \texttt{clang-format} (Google style, version-pinned) for C++,
          \texttt{rustfmt} (edition 2021, \texttt{max\_width = 100}) for Rust, and
          \texttt{gofmt} for Go. Formatting is verified by running each tool in
          \emph{check} mode and asserting a zero-diff output. Because formatting is
          entirely mechanical and has zero semantic content, any formatting violation
          is unambiguously an author oversight; the gate enforces that such oversights
          are resolved before reviewers expend attention on substantive review. All
          formatting tool versions are pinned in the repository's \texttt{.tool-versions}
          and \texttt{pyproject.toml} manifest files, ensuring deterministic formatting
          across developer workstations and CI runners alike.

    \item \textbf{Lint}: Zero warnings are permitted from \texttt{ruff} (configured
          with the full \texttt{E}, \texttt{F}, \texttt{W}, \texttt{N}, \texttt{UP},
          and \texttt{ANN} rule sets enabled), \texttt{clang-tidy} (with the
          \texttt{modernize-*}, \texttt{readability-*}, and \texttt{performance-*}
          check families active), \texttt{clippy} (invoked with
          \texttt{-D warnings -D clippy::pedantic}), and \texttt{golangci-lint}
          (running \texttt{staticcheck}, \texttt{errcheck}, \texttt{govet}, and
          \texttt{ineffassign} at minimum). The lint gate closes an entire class of
          latent defects—unused imports, shadowed variables, unchecked error returns,
          unnecessary allocations, and misused APIs—before any human reviewer is
          involved. Rule sets are reviewed and potentially extended on a quarterly
          cadence; newly activated rules are introduced with a dedicated remediation
          sprint to avoid gate failures on existing code.

    \item \textbf{Static analysis}: Zero findings at any severity level are permitted
          from the language-level static analysis tier. For Python, \texttt{mypy} is
          invoked in strict mode (\texttt{--strict --disallow-any-generics
          --warn-return-any}), requiring fully annotated signatures for all public and
          private functions, and prohibiting implicit \texttt{Any} propagation. For
          C++, a dedicated static analyser pass (Clang Static Analyzer or an equivalent
          commercial tool) is run over all translation units, checking for null
          dereferences, use-after-free, uninitialized memory reads, and integer overflow
          paths. For Rust, the borrow checker, lifetime checker, and \texttt{unsafe}
          block auditor are run as part of \texttt{cargo check --all-features}, and any
          \texttt{unsafe} block introduced since the last release must be accompanied by
          a formal safety comment documenting the invariants it relies upon. The static
          analysis gate collectively provides a machine-verified approximation of
          memory safety and type soundness across the entire polyglot codebase.

    \item \textbf{Build}: All language targets must compile without errors or warnings
          under a clean, hermetically isolated build environment. Build hermiticity is
          enforced by running all compilation steps inside a pinned container image with
          no network access and a read-only dependency cache derived from the repository's
          lockfiles (\texttt{poetry.lock}, \texttt{Cargo.lock}, \texttt{go.sum},
          \texttt{conanfile.lock}). Any implicit dependency on the host environment—an
          undeclared system library, a globally installed tool, an environment variable
          side-channel—will cause the hermetic build to fail even if the developer's
          local build succeeds. Build warnings are treated identically to errors by
          passing \texttt{-Werror} (C++), \texttt{RUSTFLAGS="-D warnings"} (Rust), and
          equivalent flags for all other targets. The resulting build artefacts are
          content-addressed and stored in the artefact registry for use by downstream
          gates, ensuring that test gates execute against exactly the binaries that
          passed compilation rather than rebuilding from source.

    \item \textbf{Unit tests}: All test categories enumerated in
          Table~\ref{tab:testing_taxonomy} must pass sequentially in the order listed.
          Test execution is isolated: each test case runs in its own subprocess with a
          freshly initialised state, a deterministic random seed derived from the test
          name, and no shared mutable global state. Tests are further partitioned into
          shards and executed in parallel across CI runner nodes, with shard assignment
          determined by a consistent-hash of the test identifier to maintain
          reproducible assignment across runs. Coverage is measured at the branch level;
          the minimum aggregate branch coverage threshold is 90\%, and any new code
          path introduced by a pull request that is not exercised by at least one test
          case causes the gate to fail. Flaky tests—those that do not produce identical
          results across three consecutive re-executions on identical inputs—are
          automatically quarantined, filed as high-priority defects, and excluded from
          gate coverage accounting until repaired.

    \item \textbf{Temporal leakage gate}: The \texttt{TemporalLeakageDetector} is
          executed against the full feature operator graph using the most recent
          historical dataset partition. The detector constructs, for each operator
          node $v$ in the directed acyclic computation graph $\mathcal{G}$, the set of
          data dependencies $\mathcal{D}(v)$, and verifies that for every dependency
          $d \in \mathcal{D}(v)$ with observation timestamp $\tau_d$ and operator
          evaluation timestamp $\tau_v$, the strict inequality $\tau_d < \tau_v$ holds.
          Any operator that consumes data with $\tau_d \geq \tau_v$—even by a single
          millisecond—is flagged as a temporal leakage violation and the gate fails
          immediately. In addition, the detector performs a statistical cross-validation
          test: model performance metrics computed on the in-sample partition are
          compared against an out-of-sample holdout, and an anomalously small
          performance gap (below a configurable threshold $\epsilon_{\text{gap}}$)
          triggers a leakage warning that is escalated to a human reviewer. This gate
          is the single most important correctness gate in the pipeline, as temporal
          leakage produces strategies that appear profitable in simulation but are
          fundamentally non-executable in live trading.

    \item \textbf{Security scan}: The security gate comprises two independent
          sub-checks that must both pass with zero critical findings. The first
          sub-check is secret scanning: the full repository history, including all
          commits reachable from the pull request branch, is scanned for high-entropy
          strings, known API key patterns, private key PEM headers, and credential
          patterns using \texttt{trufflehog} and \texttt{gitleaks}. Any detected
          credential—even one that has been revoked—causes an immediate gate failure
          and triggers an automated credential rotation workflow. The second sub-check
          is software composition analysis (SCA): all direct and transitive
          dependencies across all language ecosystems are cross-referenced against the
          National Vulnerability Database (NVD) and the GitHub Advisory Database.
          Findings at CVSS severity 7.0 or above are treated as critical and block
          the gate; findings between 4.0 and 6.9 are treated as warnings and must
          be acknowledged with a documented risk acceptance before the gate is
          considered passing. The SCA database is refreshed every six hours on CI
          runners to ensure that newly published advisories are reflected promptly.

    \item \textbf{Performance regression}: Each operator in the pipeline has a
          documented $p_{99}$ latency budget specified in its interface contract,
          expressed in microseconds. The performance regression gate executes a
          standardised benchmark suite against the build artefacts from the Build
          gate, using a dedicated, bare-metal benchmark runner with a fixed CPU
          affinity mask, disabled hyper-threading, and a pinned CPU frequency governor
          to minimise measurement variance. The $p_{99}$ latency is estimated using
          HDR histograms over $10^5$ iterations per operator, with the first $10^3$
          iterations discarded as JIT and cache warm-up. If any operator's measured
          $p_{99}$ latency exceeds its budget by more than 5\%, the gate fails and
          the pull request author receives a structured performance regression report
          identifying the offending operator, the measured versus budgeted latency,
          and the commit that introduced the regression (via binary search over the
          commit history using \texttt{git bisect} executed automatically by the CI
          system). Latency budgets are reviewed at the start of each quarter and
          may be tightened—but never relaxed without a formal architectural review—as
          operator implementations mature.

    \item \textbf{Documentation validation}: The documentation gate enforces two
          independent correctness properties over the codebase's documentation
          artefacts. First, all public interfaces—functions, methods, classes,
          modules, and configuration schema fields—must have non-empty, structurally
          valid docstrings that include at minimum a one-sentence summary, a
          \texttt{Parameters} section for every non-trivial argument, a
          \texttt{Returns} section, and a \texttt{Raises} section listing all
          documented exception types. Compliance is checked programmatically using
          \texttt{pydocstyle} (convention \texttt{numpy}) for Python and equivalent
          tooling for other languages. Second, this whitepaper and all companion
          formal specification documents must compile without errors or warnings
          using \texttt{pdflatex} or \texttt{lualatex} under a pinned \TeX{} Live
          distribution. Broken cross-references, undefined labels, missing bibliography
          entries, and overfull hboxes exceeding 5pt are all treated as warnings that
          must be resolved before the gate passes. This ensures that the formal
          documentation remains a living, compilable artefact that accurately reflects
          the implemented system rather than drifting into obsolescence.
\end{enumerate}



% ============================================================
% CHAPTER 10: OBSERVABILITY, PROVENANCE, AND AUDIT
% ============================================================
\chapter{Observability, Provenance, and Audit}

\section{The Reconstruction Requirement}

For any decision produced by the platform at time $t$, it must be possible
to reconstruct, post hoc, the complete causal chain:

\begin{enumerate}
    \item Which raw data were ingested, with what provenance attributes.
    \item Which normalisation and feature operators ran, with what versions.
    \item What regime posterior distribution was produced, by which model.
    \item What signals were generated, with what confidence scores.
    \item What risk evaluation was performed, with what outcome.
    \item What order was submitted or paper-executed.
    \item What fill was received (or simulated).
    \item What the resulting portfolio state was.
\end{enumerate}

The reconstruction requirement is not merely a convenience for post-mortem debugging—it is a foundational epistemic guarantee that the platform makes to its operators, regulators, and auditors. In quantitative trading systems operating at high frequency or across multiple asset classes simultaneously, the causal pathway from a raw market datum to an executed order can traverse dozens of stateful operators, probabilistic models, and asynchronous message queues. Without a rigorous and formally specified reconstruction mechanism, any post-hoc attribution of a trading outcome to a specific model decision or data artefact becomes speculative at best and legally indefensible at worst. The platform therefore treats full causal reconstructibility as a hard invariant that must be preserved under all operating conditions, including partial node failure, network partition, operator restarts, and schema migrations.

Formally, let $\mathcal{D} = \{d_1, d_2, \ldots, d_n\}$ denote the set of raw data records ingested during a session window $[t_0, t_1)$. Each record $d_i$ carries a provenance tuple $\pi(d_i) = (\texttt{source\_id}, \texttt{feed\_version}, \texttt{ingestion\_timestamp}, \texttt{checksum})$ that uniquely identifies its origin and content integrity. The platform guarantees that for every downstream artefact $a$ produced by the pipeline, there exists a computable provenance function $\Pi(a) \subseteq \mathcal{D}$ that enumerates the minimal set of raw records from which $a$ was derived. This guarantee extends transitively through all intermediate operators: if operator $O_k$ consumes artefact $a_{k-1}$ and produces $a_k$, then $\Pi(a_k) = \bigcup_{a \in \mathrm{inputs}(O_k)} \Pi(a)$, and the operator's version identifier and parameter snapshot are recorded alongside the output so that the transformation itself can be replayed deterministically.

The normalisation and feature computation layer amplifies the complexity of this provenance graph considerably. A single feature vector presented to a regime classification model may depend on rolling window statistics computed over hundreds of historical bars, each of which was itself the product of corporate action adjustments, timezone normalisation, outlier truncation, and interpolation across missing values. The platform addresses this by requiring every \texttt{FeatureOperator} to emit a \texttt{FeatureLineageRecord} alongside each output batch, containing the operator class name, its semantic version, the hash of its frozen parameter configuration, the closed interval of input timestamps consumed, and the list of \texttt{correlation\_id} values for all upstream records that contributed to the computation. These lineage records are stored in a dedicated append-only lineage store, indexed by both \texttt{correlation\_id} and output feature name, enabling sub-second retrieval of the full dependency graph for any feature at any historical timestamp.

The regime posterior distribution produced by the Hidden Markov Model or Bayesian switching model introduces a probabilistic dimension to the causal chain that requires careful handling. Because the posterior $P(s_t \mid \mathbf{x}_{1:t})$ is computed by a sequential inference algorithm whose state depends on the entire history of observations $\mathbf{x}_{1:t}$, reconstruction requires not only the current observation $\mathbf{x}_t$ but also the model's internal belief state at time $t-1$. The platform therefore snapshots the full belief state vector—represented as a probability simplex over the $K$ latent regime states—at every inference step and persists it to the audit store with a pointer to the model artefact (identified by its content-addressed hash in the model registry). This ensures that the exact posterior can be reproduced by re-running forward inference from the snapshotted belief state, even if the upstream observation history has been compacted or archived.

Signal generation produces one or more typed signal objects, each carrying a confidence score derived from the regime posterior, the feature values, and the strategy-specific decision function. The reconstruction requirement mandates that the exact floating-point values of all inputs to the confidence scoring function be recorded, since even single-bit differences in intermediate computations can propagate through nonlinear transformations to produce materially different signal outcomes. To satisfy this requirement without incurring prohibitive storage overhead, the platform employs a selective snapshotting strategy: all inputs are hashed using a collision-resistant hash function (SHA-256), and the full input tensor is stored only when the resulting signal crosses a materiality threshold defined in the strategy configuration. For signals below the threshold, the hash alone is retained, enabling integrity verification without full replay capability. Above the threshold—that is, for any signal that influences an order decision—the complete input snapshot is mandatory and its absence is treated as a fatal audit violation.

Risk evaluation applies a sequence of constraint checks and portfolio-level exposure computations to each candidate signal before authorising order submission. The reconstruction requirement for this stage demands that the exact portfolio state used as input to the risk engine—including all open positions, pending orders, margin utilisation, and drawdown watermarks—be recoverable for the timestamp at which the evaluation occurred. The platform satisfies this by maintaining a time-indexed portfolio state journal: every mutation to the portfolio state (position open, position close, fill receipt, margin call) is recorded as an immutable event in an event-sourced ledger, and the portfolio state at any past timestamp can be reconstructed by replaying the event stream up to that point. The risk evaluation record then stores only the timestamp and the hash of the reconstructed state, keeping storage costs proportional to the number of state-mutating events rather than the number of risk evaluations.

Order submission and fill receipt are the two stages at which the platform's internal state interacts with the external world, and it is precisely at these boundaries that provenance chains are most vulnerable to breakage. Network failures, exchange-side rejections, partial fills, and latency-induced race conditions can all create discrepancies between the order the platform intended to submit and the order that was actually acknowledged by the broker. The platform mitigates this by treating every outbound order as a tentative event that transitions through a formally specified state machine: \texttt{PENDING\_SUBMIT} $\to$ \texttt{SUBMITTED} $\to$ \texttt{ACKNOWLEDGED} $\to$ \texttt{PARTIALLY\_FILLED} $\to$ \texttt{FILLED} (or \texttt{REJECTED} / \texttt{CANCELLED}). Each state transition is logged as an immutable audit event with a monotonically increasing sequence number, and the platform will not advance to a subsequent pipeline stage for a given \texttt{correlation\_id} until the order state machine has reached a terminal state or an explicit timeout has been recorded and escalated.

Every object passing through the pipeline carries a \texttt{correlation\_id}—a UUID (version 4, randomly generated at ingestion time) that is propagated without modification through all transformations, fan-outs, and aggregations. When a single upstream event gives rise to multiple downstream artefacts (as occurs when a market data tick triggers feature recomputation for multiple strategies simultaneously), each derived artefact retains the original \texttt{correlation\_id} and additionally records a \texttt{parent\_correlation\_id} pointer forming a directed acyclic graph of causal relationships. This DAG is materialised in the audit store and can be queried using a graph traversal interface, enabling operators to answer questions of the form: ``For all orders submitted between $t_a$ and $t_b$ that resulted in a loss exceeding $\delta$ dollars, which raw data sources contributed to the signals that generated those orders, and did any of those sources share a common upstream feed that experienced anomalous latency during that interval?'' The ability to answer such queries efficiently and with formal correctness guarantees is what elevates the platform's observability infrastructure from a passive logging mechanism to an active instrument of risk governance.

\section{Structured Logging Requirements}

The platform uses structured (JSON) logs exclusively. Every log record must
include:

\begin{table}[H]
\centering
\caption{Required fields in every structured log record.}
\label{tab:log_fields}
\small
\begin{tabularx}{\textwidth}{p{4.2cm} p{10.3cm}}
\toprule
\textbf{Field} & \textbf{Description} \\
\midrule
\texttt{timestamp} & UTC ISO-8601 timestamp with millisecond precision. \\
\texttt{level} & Log level: \texttt{DEBUG}, \texttt{INFO}, \texttt{WARN},
    \texttt{ERROR}, \texttt{CRITICAL}. \\
\texttt{operator\_id} & Unique identifier of the operator emitting the record. \\
\texttt{version} & Operator version string. \\
\texttt{correlation\_id} & UUID linking this record to the originating
    pipeline invocation. \\
\texttt{run\_id} & Identifier of the simulation run or live session. \\
\texttt{event\_type} & Structured event type from the event taxonomy. \\
\texttt{payload} & Event-specific structured data. \\
\bottomrule
\end{tabularx}
\end{table}

\section{Audit Record Lifecycle}

The \texttt{ValidationOperator} at each pipeline stage emits an
\texttt{AuditRecord} that is persisted to an append-only audit log.
The audit log is tamper-evident: each record includes a hash of the
previous record, forming a cryptographic chain. Any post-hoc modification
of the audit log is detectable.

% ============================================================
% CHAPTER 11: DEPLOYMENT AND OPERATIONAL CONSIDERATIONS
% ============================================================
\chapter{Deployment and Operational Considerations}

\section{Environment Separation}

The platform enforces a rigorous, multi-layered separation between four distinct operational environments: \emph{research}, \emph{simulation}, \emph{paper trading}, and \emph{live production}. This separation is not merely a matter of convention or documentation; it is mechanically enforced at every level of the software stack, from the build system through to runtime assertions embedded in operator implementations. The rationale for this strictness stems from the asymmetric consequence model of trading systems: a misrouted order in production that was intended for a simulation environment can result in real financial loss, regulatory exposure, and reputational damage that far exceeds the cost of imposing robust environmental controls.

Environment-specific configuration is loaded exclusively from environment variables and validated exhaustively against a typed schema at process startup. Any missing or malformed variable causes the process to abort with a structured error report before any operator is instantiated. This fail-fast posture prevents a partially-configured process from reaching a state where it might silently behave as a different environment. The schema validation is performed by a dedicated \texttt{EnvironmentBootstrapper} component that runs prior to dependency injection and operator graph construction, ensuring that no downstream component ever observes an inconsistent configuration state.

Cross-contamination between environments is prevented by a defence-in-depth strategy comprising three independent enforcement layers. First, each environment maintains entirely separate configuration files, secrets stores, broker endpoint registries, and database connection pools. No file path, network socket, or storage volume is shared across environment boundaries; infrastructure-as-code tooling enforces this at the provisioning layer, and access control lists at the cloud provider level further restrict cross-environment network reachability. Second, the C++ and Rust compilation units that implement order routing and execution submission are guarded by compile-time feature flags (\texttt{\#ifdef LIVE\_EXECUTION} in C++ and Cargo feature gates in Rust) that unconditionally disable all live order-submission code paths in builds not explicitly tagged for production. This means that a simulation binary is physically incapable of submitting a live order, regardless of configuration, because the relevant code is absent from the compiled artefact. Third, every concrete implementation of \texttt{ExecutionOperator} contains a runtime assertion, evaluated on each invocation, that compares the active environment token embedded in the operator context against the environment token baked into the binary at compile time. If the two tokens disagree, the operator raises a \texttt{EnvironmentMismatchError} and halts the pipeline immediately, writing a \texttt{CRITICAL}-level audit record before terminating.

\begin{enumerate}
    \item Separate configuration files and secrets stores per environment, with infrastructure-level access controls enforcing the boundary.
    \item Compile-time flags in C++ and Rust that unconditionally exclude live execution code paths from non-production build artefacts.
    \item Runtime assertion in all \texttt{ExecutionOperator} implementations that the active environment token matches the compiled-in environment token, with immediate pipeline halt on mismatch.
\end{enumerate}

The research environment additionally enforces read-only access to all market data stores, preventing accidental mutation of historical datasets that would corrupt the reproducibility of backtests. The simulation environment is granted write access only to ephemeral namespaces within the time-series database, and all records written therein are tagged with a \texttt{sim:} prefix that prevents them from being ingested by live data pipelines. Paper trading connects to broker sandbox APIs that are contractually guaranteed by the broker to never route orders to an exchange; the platform additionally validates, on each order acknowledgement, that the returned order identifier conforms to the sandbox identifier format defined in the broker's API specification.

\section{Secret Management}

The platform enforces an absolute prohibition on the presence of credentials, API keys, broker authentication tokens, database passwords, cryptographic private keys, or any other secret material in source code, configuration files, build artefacts, log output, or inter-process communication payloads. This prohibition is not advisory: the CI/CD pipeline incorporates a secret-scanning gate, powered by a combination of entropy-based heuristics and a curated pattern library covering the credential formats of all brokers, cloud providers, and data vendors used by the platform. Any commit that triggers a detection is automatically rejected at the pull-request level, and the offending developer is notified with a remediation guide. The pattern library is versioned alongside the codebase and updated whenever a new credential format is introduced into the dependency graph.

At runtime, secrets are sourced exclusively from two mechanisms. For local development and CI workers, secrets are injected as environment variables by the developer's shell profile or by the CI secret store (currently HashiCorp Vault, accessed via the Vault Agent Sidecar injector pattern). For production deployments, secrets are fetched at process startup by the \texttt{SecretsBootstrapper} component, which authenticates to the secrets manager using a short-lived, automatically-rotated machine identity credential issued by the cloud provider's identity and access management layer. The fetched secrets are held in memory in a \texttt{SecretString} type whose destructor overwrites the backing buffer with zeros before deallocation, reducing the window during which secret material is resident in a potentially-swappable memory page. The platform also disables core dumps in production, and the process memory lock flag (\texttt{mlockall} on Linux) is set to prevent secret-bearing pages from being swapped to disk.

Secret rotation is handled without process restart: the \texttt{SecretsBootstrapper} subscribes to rotation events from the secrets manager and, upon receiving a rotation notification, fetches the new secret version and atomically replaces the in-memory reference using a reader-writer lock. Operators that hold a reference to a secret value access it through a \texttt{SecretRef} handle that dereferences through the lock on each use, ensuring that the first invocation after a rotation always observes the new value. All secret access events—including the initial fetch, each rotation, and each dereference—are recorded in the audit log with the secret's logical name and version identifier but never with the secret value itself.

\section{Configuration Philosophy}

The platform draws a principled and strictly enforced boundary between \emph{behaviour}, which resides exclusively in source code, and \emph{parameters}, which reside exclusively in externalised configuration files. All changeable parameters—including risk limits per instrument and portfolio, model hyperparameters for each operator variant, instrument universe definitions, data source endpoint addresses, simulation time horizons, fee schedules, slippage model coefficients, and execution algorithm tuning constants—are expressed in TOML or YAML configuration files that are loaded at startup and never modified at runtime. The choice of TOML for operator-level configuration and YAML for infrastructure-level configuration reflects the differing readability and schema-expressiveness requirements of each layer: TOML's explicit typing and prohibition of implicit nulls reduce the risk of misconfiguration in numerically-sensitive operator parameters, while YAML's support for anchors and aliases reduces repetition in the multi-environment infrastructure configurations.

Every configuration file is validated at startup against a JSON Schema document that precisely specifies the type, range, unit, and inter-field constraint for each parameter. For example, the risk configuration schema asserts that \texttt{max\_position\_fraction} must be a floating-point number in the open interval $(0, 1)$, that \texttt{max\_drawdown\_halt\_threshold} must be strictly less than \texttt{max\_drawdown\_warn\_threshold}, and that the set of instruments listed in the universe must be a non-empty subset of the instruments for which fee schedules are defined. Schema violations produce structured error messages that identify the offending field, the violated constraint, and the observed value, enabling rapid diagnosis without requiring the operator to inspect source code. The schema documents are themselves versioned and their compatibility with the running codebase is verified by a startup check that compares schema version hashes.

Configuration files contain no executable logic whatsoever. Expressions, templating directives, scripting syntax, and embedded code are not permitted. This constraint is enforced by the configuration loader, which uses a strict-mode parser that rejects any input not conformable to the grammar of a pure data document. The motivation for this constraint is the reproducibility and auditability of system behaviour: a configuration file that contains executable logic can produce different effective configurations depending on the environment in which it is evaluated, undermining the guarantee that a known configuration file corresponds to a known, fixed set of parameter values. By restricting configuration to pure data, the platform ensures that a configuration file can be fully characterised by a content hash, and that two processes started with identical configuration file hashes and identical binary hashes will exhibit identical parameterisation.

\section{Dependency Discipline}

The platform treats the introduction of a new third-party dependency as a significant architectural decision requiring explicit justification, not a trivial implementation convenience. This posture reflects the long-term costs that accrue from unchecked dependency growth: increased attack surface, supply-chain risk, build time inflation, licence compliance burden, and the cognitive load imposed on future maintainers who must understand each dependency's behaviour and failure modes. Every proposed new dependency must pass a six-criterion gate before it is approved for inclusion in the dependency manifest.

The first criterion requires a documented demonstration that the C++ or Rust standard library does not already provide an adequate solution. Standard library facilities are preferred unconditionally because they carry no additional supply-chain risk, are subject to the most rigorous review of any code in the ecosystem, and impose no licence obligations beyond those already accepted. The second criterion requires a review of all existing approved dependencies to confirm that none of them already expose the required functionality, either directly or via a minor extension of their public interface. The third criterion requires evidence that the candidate dependency is actively maintained, defined as at least one substantive commit to its primary repository within the preceding twelve months and a responsive maintainer team with a documented process for accepting security disclosures. The licence must be compatible with the platform's own licence and with the licences of all existing dependencies, as verified by an automated licence-compatibility checker integrated into the dependency review workflow. The fourth criterion requires a scan of all known vulnerability databases (NVD, OSV, and the GitHub Advisory Database) confirming that the dependency has no open critical or high-severity vulnerabilities; any dependency with an open vulnerability of this severity is ineligible until a patched version is available and the platform has migrated to it. The fifth criterion requires a benchmark demonstrating that the runtime and compile-time costs introduced by the dependency are justified by the benefit it delivers, with specific reference to the latency budgets defined in Section~\ref{sec:performance_budgets}. The sixth criterion requires an assessment of cross-platform implications, including compatibility with all supported operating systems and processor architectures, and any platform-specific behaviour differences that could cause divergence between development and production environments.

\begin{enumerate}
    \item The standard library does not already solve the problem, as demonstrated by explicit review of the relevant standard library documentation.
    \item No existing approved dependency already solves the problem, as demonstrated by review of the current dependency manifest and its exposed APIs.
    \item The dependency is actively maintained (substantive commit activity within twelve months), has a responsive security disclosure process, and carries a licence compatible with all existing dependencies and the platform's own licence.
    \item The dependency has no open critical or high-severity vulnerabilities in NVD, OSV, or the GitHub Advisory Database as of the date of the review.
    \item The runtime and compile-time costs introduced by the dependency are quantified by benchmark and are justified by the benefit, with explicit reference to the operator latency budgets.
    \item Cross-platform implications have been assessed, including compatibility with all supported operating systems and architectures, and any platform-specific behavioural differences have been documented.
\end{enumerate}

Approved dependencies are pinned to exact versions in the manifest (using Cargo's \texttt{Cargo.lock} for Rust and a pip-compiled \texttt{requirements.txt} with hash verification for Python tooling), and automatic dependency update pull requests generated by Dependabot are gated on the same six-criterion review before being merged. This ensures that even routine patch-level updates receive security and compatibility scrutiny before reaching the main branch.

\section{Performance Budgets}
\label{sec:performance_budgets}

The platform defines explicit, per-category latency budgets for every class of operator in the pipeline, expressed as $p_{99}$ wall-clock latency targets measured under representative production load on the reference hardware configuration. The $p_{99}$ metric is chosen rather than the mean or median because the pipeline is sensitive to tail latency: a single slow operator invocation at the $p_{99}$ or beyond will delay the entire pipeline tick, potentially causing the downstream \texttt{ExecutionOperator} to submit an order stale by one or more market data updates. The budgets are not aspirational targets but hard service-level objectives: any operator implementation that exceeds its budget at the $p_{99}$ percentile in the pre-deployment benchmark suite causes the CI pipeline to fail and blocks the release.

The budget for the \texttt{DataOperator} category is set at 50 milliseconds at $p_{99}$, which is sufficient to accommodate a single round-trip network request to a co-located market data feed under normal network conditions, including the serialisation and deserialisation overhead of the chosen wire format. The \texttt{FeatureOperator} budget of 5 milliseconds reflects the fact that feature computation is performed entirely in-process on already-fetched data, and that the most computationally intensive features—rolling exponentially-weighted covariance matrices over a 252-bar lookback—have been benchmarked to complete in under 3 milliseconds on the reference hardware. The \texttt{MarketRegimeOperator} budget of 20 milliseconds accounts for the ensemble aggregation step in which outputs from multiple regime-detection sub-models (Hidden Markov Models, volatility-scaled momentum indicators, and cross-asset correlation monitors) are combined via a weighted soft-voting scheme whose weights are themselves updated on a rolling basis. The \texttt{SignalOperator} budget of 2 milliseconds per signal is the tightest non-execution budget in the pipeline and reflects the requirement that signal evaluation be reducible to a small number of vectorised floating-point operations on pre-computed feature values. Any signal design that requires iterative numerical solving or Monte Carlo sampling is categorically incompatible with this budget and must be reformulated or relegated to a lower-frequency pipeline tier. The \texttt{RiskOperator} budget of 10 milliseconds includes the time required to evaluate Conditional Value-at-Risk at the 95\% confidence level using a historical simulation over a 500-scenario window, as well as the evaluation of all per-instrument and portfolio-level hard limits. The \texttt{PortfolioOperator} budget of 15 milliseconds encompasses the convex optimisation solve for mean-variance or risk-parity portfolio construction, using a warm-started interior-point solver whose initial point is the previous period's optimal allocation; empirical benchmarks show that the solver converges in under 12 milliseconds for universes of up to 50 instruments. The \texttt{ExecutionSimulator} budget of 1 millisecond applies exclusively to the paper-trading fill simulation path, which involves no network I/O and consists solely of in-memory order-book traversal and fill model evaluation.

\begin{table}[H]
\centering
\caption{Operator latency budgets by category.}
\label{tab:latency_budgets}
\small
\begin{tabularx}{\textwidth}{p{4.5cm} p{4.5cm} p{5.5cm}}
\toprule
\textbf{Operator Category} & \textbf{$p_{99}$ Latency Budget} & \textbf{Notes} \\
\midrule
\texttt{DataOperator} & 50 ms & Includes network round trip to co-located data feed, plus wire-format deserialisation. \\
\texttt{FeatureOperator} & 5 ms & In-process computation only; rolling EW covariance benchmarked at $<$3 ms on reference hardware. \\
\texttt{MarketRegimeOperator} & 20 ms & Includes HMM inference, volatility-scaled momentum evaluation, and ensemble soft-vote aggregation. \\
\texttt{SignalOperator} & 2 ms & Per signal, excluding regime; requires signal reducibility to vectorised operations on pre-computed features. \\
\texttt{RiskOperator} & 10 ms & Includes CVaR at 95\% over 500-scenario historical simulation window and all hard-limit evaluations. \\
\texttt{PortfolioOperator} & 15 ms & Includes warm-started interior-point convex optimisation solve for universes up to 50 instruments. \\
\texttt{ExecutionSimulator} & 1 ms & Paper fill only; no network I/O; in-memory order-book traversal and fill model evaluation. \\
\bottomrule
\end{tabularx}
\end{table}

% ============================================================
% CHAPTER 12: SUMMARY OF KEY RESULTS AND FORMULAS
% ============================================================
\chapter{Summary of Key Results and Formulas}

\begin{table}[H]
\centering
\caption{Summary of core mathematical formulas used across the platform (Part I).}
\label{tab:formulas_1}
\small
\begin{tabularx}{\textwidth}{p{3.8cm} p{11.0cm}}
\toprule
\textbf{Formula} & \textbf{Expression and Context} \\
\midrule
Kelly fraction &
    $f^* = p - q/b$, where $p$ is win probability, $q=1-p$, $b$ is payoff ratio.
    Used by \texttt{PositionSizingOperator}. \\
\midrule
Confidence-weighted signal &
    $S = \sum_i c_i s_i / \sum_i c_i$, with $S \in [-1,1]$.
    Used by \texttt{ConfidenceAggregationOperator}. \\
\midrule
HMM belief update &
    Equation~\eqref{eq:hmm_filter}: Bayesian filter update for regime posterior.
    Used by \texttt{MarketRegimeOperator}. \\
\midrule
Regime entropy &
    $H_t = -\sum_k \beta_t(k) \ln \beta_t(k)$, normalised as $\hat{U}_t = H_t/\ln K$.
    Used to gate conservative sizing. \\
\midrule
Minimum-variance hedge ratio &
    $h^* = \rho_{VF} \sigma_V / \sigma_F$, equation~\eqref{eq:hedge_ratio}.
    Used by \texttt{HedgeOptimizationOperator}. \\
\bottomrule
\end{tabularx}
\end{table}

\begin{table}[H]
\centering
\caption{Summary of core mathematical formulas used across the platform (Part II).}
\label{tab:formulas_2}
\small
\begin{tabularx}{\textwidth}{p{3.8cm} p{11.0cm}}
\toprule
\textbf{Formula} & \textbf{Expression and Context} \\
\midrule
CVaR (LP formulation) &
    Equation~\eqref{eq:cvar_lp}: $\CVaR_\alpha = \min_z \{z + \frac{1}{(1-\alpha)T}\sum_t \max(\ell_t - z, 0)\}$.
    Used by \texttt{RiskOperator}. \\
\midrule
Risk-adjusted objective &
    $\mathcal{J}(\pi) = \mathbb{E}[\sum_t \gamma^t r_t] - \lambda_1 \CVaR_\alpha - \lambda_2 \text{MaxDD}$,
    equation~\eqref{eq:objective}. \\
\midrule
Cross-domain allocation &
    $\boldsymbol{\alpha}^* = \arg\max_{\alpha \in \Delta^3}[\alpha^\top \mu^{(r)} - \lambda \alpha^\top \Sigma^{(r)} \alpha]$,
    equation~\eqref{eq:cross_domain_alloc}. \\
\midrule
Market impact &
    $\text{Impact} = \gamma \sigma \sqrt{V/\text{ADV}}$.
    Used by \texttt{TransactionCostOperator}. \\
\midrule
Temporal validity invariant &
    $\forall t: \text{Decision}_t(\pi) \in \sigma(\mathcal{I}_t)$.
    Enforced by \texttt{TemporalLeakageDetector} CI gate. \\
\bottomrule
\end{tabularx}
\end{table}

% ============================================================
% REFERENCES
% ============================================================
\chapter*{References}
\addcontentsline{toc}{chapter}{References}

\begin{enumerate}
    \item Markowitz, H. (1952). Portfolio Selection. \textit{Journal of Finance}, 7(1), 77--91.
    \item Kelly, J.\,L. (1956). A New Interpretation of Information Rate.
          \textit{Bell System Technical Journal}, 35(4), 917--926.
    \item Sharpe, W.\,F. (1966). Mutual Fund Performance.
          \textit{Journal of Business}, 39(1), 119--138.
    \item Sortino, F.\,A., \& van der Meer, R. (1991). Downside Risk.
          \textit{Journal of Portfolio Management}, 17(4), 27--31.
    \item Rockafellar, R.\,T., \& Uryasev, S. (2000). Optimisation of
          Conditional Value-at-Risk. \textit{Journal of Risk}, 2(3), 21--41.
    \item Hamilton, J.\,D. (1989). A New Approach to the Economic Analysis
          of Nonstationary Time Series and the Business Cycle.
          \textit{Econometrica}, 57(2), 357--384.
    \item Ederington, L.\,H. (1979). The Hedging Performance of the New
          Futures Markets. \textit{Journal of Finance}, 34(1), 157--170.
    \item Lopez de Prado, M. (2018). \textit{Advances in Financial Machine Learning}.
          Wiley.
    \item Pardo, R. (2008). \textit{The Evaluation and Optimisation of Trading
          Strategies} (2nd ed.). Wiley.
    \item Almgren, R., \& Chriss, N. (2001). Optimal Execution of Portfolio
          Transactions. \textit{Journal of Risk}, 3(2), 5--39.
    \item Ben-Tal, A., \& Nemirovski, A. (1998). Robust Convex Optimisation.
          \textit{Mathematics of Operations Research}, 23(4), 769--805.
    \item Shannon, C.\,E. (1948). A Mathematical Theory of Communication.
          \textit{Bell System Technical Journal}, 27(3), 379--423.
\end{enumerate}

% ============================================================
% APPENDIX: DECISION RECONCILIATION MATRIX
% ============================================================
\appendix
\chapter{Decision Reconciliation Matrix}

The following table maps every major architectural decision to its
justification, alternatives considered, and the conditions under which
the decision should be revisited.

\begin{table}[H]
\centering
\caption{Architectural decision reconciliation matrix (Part I).}
\label{tab:decisions_1}
\small
\begin{tabularx}{\textwidth}{p{3.5cm} p{4.5cm} p{3.5cm} p{3.0cm}}
\toprule
\textbf{Decision} & \textbf{Justification} & \textbf{Alternatives Considered} & \textbf{Revisit Trigger} \\
\midrule
Operator composition over monolithic model &
    Small typed units are independently testable and replaceable
    without full system rewrite &
    Single-model neural network end-to-end &
    If composition overhead exceeds 10\% of total latency budget \\
\midrule
Python as research language &
    Richest quantitative ecosystem; fastest strategy iteration &
    Julia, R &
    If Python runtime becomes a bottleneck for non-numerical code \\
\midrule
Rust for execution engine &
    Memory safety and deterministic latency without GC pauses &
    C++, Go &
    If Rust ecosystem lacks required financial library support \\
\midrule
Protocol Buffers for shared schemas &
    Mature, language-neutral, generates bindings for all four languages &
    Apache Arrow, FlatBuffers, JSON Schema &
    If schema evolution friction becomes unacceptable \\
\midrule
Simulation-first deployment &
    Prevents accidental live trading; forces validation before exposure &
    Direct live deployment with kill switch &
    Never; live trading always requires explicit deliberate enablement \\
\bottomrule
\end{tabularx}
\end{table}

\begin{table}[H]
\centering
\caption{Architectural decision reconciliation matrix (Part II).}
\label{tab:decisions_2}
\small
\begin{tabularx}{\textwidth}{p{3.5cm} p{4.5cm} p{3.5cm} p{3.0cm}}
\toprule
\textbf{Decision} & \textbf{Justification} & \textbf{Alternatives Considered} & \textbf{Revisit Trigger} \\
\midrule
HMM for regime detection &
    Well-understood Bayesian framework with calibrated uncertainty;
    produces probability distributions not point labels &
    K-means clustering, threshold rules, deep learning classifier &
    If HMM posterior calibration degrades significantly in live data \\
\midrule
Kelly criterion with fractional scaling &
    Theoretically optimal for log-wealth maximisation; fraction avoids
    estimation-error sensitivity &
    Fixed fractional, equal weighting &
    If win probability and payoff ratio cannot be estimated reliably \\
\midrule
CVaR as primary tail risk metric &
    Coherent risk measure; linear programming formulation is tractable;
    better than VaR for tail events &
    VaR, maximum drawdown alone &
    Never; CVaR is strictly superior to VaR and is retained \\
\midrule
Append-only audit log with hash chain &
    Tamper-evident; any modification detectable post hoc &
    Mutable relational database &
    If storage costs become prohibitive; consider archival tiering \\
\midrule
Temporal leakage as CI gate &
    Makes correctness structural, not advisory; zero tolerance &
    Code review only &
    Never; temporal correctness is a hard requirement \\
\bottomrule
\end{tabularx}
\end{table}



% ============================================================
% REFERENCES
% ============================================================
\chapter*{References}
\addcontentsline{toc}{chapter}{References}

\begin{thebibliography}{99}

\bibitem{kelly1956}
Kelly, J.~L. (1956).
A new interpretation of information rate.
\textit{Bell System Technical Journal}, 35(4), 917--926.

\bibitem{rockafellar2000}
Rockafellar, R.~T., \& Uryasev, S. (2000).
Optimization of conditional value-at-risk.
\textit{Journal of Risk}, 2(3), 21--41.

\bibitem{hamilton1989}
Hamilton, J.~D. (1989).
A new approach to the economic analysis of nonstationary time series and the business cycle.
\textit{Econometrica}, 57(2), 357--384.

\bibitem{rabiner1989}
Rabiner, L.~R. (1989).
A tutorial on hidden Markov models and selected applications in speech recognition.
\textit{Proceedings of the IEEE}, 77(2), 257--286.

\bibitem{markowitz1952}
Markowitz, H. (1952).
Portfolio selection.
\textit{The Journal of Finance}, 7(1), 77--91.

\bibitem{blacklitterman1992}
Black, F., \& Litterman, R. (1992).
Global portfolio optimization.
\textit{Financial Analysts Journal}, 48(5), 28--43.

\bibitem{Lopez2018}
Lopez de Prado, M. (2018).
\textit{Advances in Financial Machine Learning}.
John Wiley \& Sons.

\bibitem{avellaneda2010}
Avellaneda, M., \& Lee, J.-H. (2010).
Statistical arbitrage in the US equities market.
\textit{Quantitative Finance}, 10(7), 761--782.

\bibitem{fama1970}
Fama, E.~F. (1970).
Efficient capital markets: A review of theory and empirical work.
\textit{The Journal of Finance}, 25(2), 383--417.

\bibitem{grinold1999}
Grinold, R.~C., \& Kahn, R.~N. (1999).
\textit{Active Portfolio Management: A Quantitative Approach for Producing Superior Returns and Controlling Risk} (2nd ed.).
McGraw-Hill.

\bibitem{engle1982}
Engle, R.~F. (1982).
Autoregressive conditional heteroscedasticity with estimates of the variance of United Kingdom inflation.
\textit{Econometrica}, 50(4), 987--1007.

\bibitem{bollerslev1986}
Bollerslev, T. (1986).
Generalized autoregressive conditional heteroskedasticity.
\textit{Journal of Econometrics}, 31(3), 307--327.

\bibitem{litterman1991}
Litterman, R., \& Scheinkman, J. (1991).
Common factors affecting bond returns.
\textit{The Journal of Fixed Income}, 1(1), 54--61.

\bibitem{carhart1997}
Carhart, M.~M. (1997).
On persistence in mutual fund performance.
\textit{The Journal of Finance}, 52(1), 57--82.

\bibitem{fama1993}
Fama, E.~F., \& French, K.~R. (1993).
Common risk factors in the returns on stocks and bonds.
\textit{Journal of Financial Economics}, 33(1), 3--56.

\bibitem{bishop2006}
Bishop, C.~M. (2006).
\textit{Pattern Recognition and Machine Learning}.
Springer.

\bibitem{goodfellow2016}
Goodfellow, I., Bengio, Y., \& Courville, A. (2016).
\textit{Deep Learning}.
MIT Press.

\bibitem{sutton2018}
Sutton, R.~S., \& Barto, A.~G. (2018).
\textit{Reinforcement Learning: An Introduction} (2nd ed.).
MIT Press.

\bibitem{nakamoto2008}
Nakamoto, S. (2008).
Bitcoin: A peer-to-peer electronic cash system.
\textit{Cryptography Mailing List at metzdowd.com}.

\bibitem{burns2021}
Burns, B., Grant, B., Oppenheimer, D., Brewer, E., \& Wilkes, J. (2016).
Borg, Omega, and Kubernetes.
\textit{ACM Queue}, 14(1), 70--93.

\bibitem{dean2004}
Dean, J., \& Ghemawat, S. (2004).
MapReduce: Simplified data processing on large clusters.
\textit{Proceedings of the 6th USENIX Symposium on Operating Systems Design and Implementation (OSDI)}, 137--150.

\bibitem{zaharia2016}
Zaharia, M., Xin, R.~S., Wendell, P., Das, T., Armbrust, M., Dave, A., et~al. (2016).
Apache Spark: A unified engine for big data processing.
\textit{Communications of the ACM}, 59(11), 56--65.

\bibitem{kleppmann2017}
Kleppmann, M. (2017).
\textit{Designing Data-Intensive Applications}.
O'Reilly Media.

\bibitem{fowler2002}
Fowler, M. (2002).
\textit{Patterns of Enterprise Application Architecture}.
Addison-Wesley.

\end{thebibliography}



% ============================================================
% CHANGELOG
% ============================================================
\chapter*{Changelog}
\addcontentsline{toc}{chapter}{Changelog}

\begin{table}[H]
\centering
\caption{Document version history.}
\label{tab:changelog}
\small
\begin{tabularx}{\textwidth}{p{2.8cm} p{2.8cm} p{2.8cm} p{6.0cm}}
\toprule
\textbf{Version} & \textbf{Date} & \textbf{Author} & \textbf{Description} \\
\midrule
2026.1.0.0 & 2025-08-18 & Hadrian Hu &
    Initial creation. Full high-level architecture whitepaper covering
    four-domain model, operator composition framework, data pipeline,
    regime engine, risk management, polyglot architecture, testing
    taxonomy, observability, and deployment. \\
\bottomrule
\end{tabularx}
\end{table}

\end{document}
